% PNAS main text, draft v2 (2026-09-29). Built from main.tex per pnas-plan-2026-09-29.md.
% Layout: Aydin's standard template (amsart, 11pt, onehalfspacing, 4 cm margins), same as
% main.tex; PNAS will reflow into its own template. Content organization follows PNAS:
% Significance, untitled introduction, Results, Discussion, Materials and Methods; SI
% Appendix separate. Citations natbib author-year here; PNAS uses numbered.
\documentclass[article, 11pt, reqno]{amsart}
\usepackage[utf8]{inputenc}
\usepackage[english]{babel}
\usepackage{mathtools, amsmath, amsthm, amsfonts, amssymb} % Math packages
\usepackage{bm}

\usepackage[hidelinks]{hyperref} % for hyper references
\usepackage{cleveref}
\crefname{section}{§\hspace{-0.1cm}}{§§}
\Crefname{section}{§}{§§}

%--------------------------------------
% For controlling page margins
\usepackage[top=3cm, bottom=3cm, left=4cm, right=4cm]{geometry}
% For bibliography and citation style
\usepackage{csquotes}
\usepackage{natbib}
\setcitestyle{aysep={}}
\usepackage[toc,page]{appendix} % For linking to the appendix
%--------------------------------------
% text spacing, alignment, and formatting
\usepackage{setspace}
\usepackage[multiple,hang,flushmargin]{footmisc} % Remove indent from footnotes, allow multiple footnotes
\usepackage[foot]{amsaddr}
\setlength{\skip\footins}{.3cm} % Control spacing between footnotes and body text
\onehalfspacing
\frenchspacing
\raggedbottom
%--------------------------------------
% aesthetic tweaks for quoting, figures, and tables
\usepackage{upquote} % For `' quote formatting
\usepackage{graphicx} % For importing images
\usepackage{booktabs} % Improves table options & aesthetics
\usepackage{caption} % For captions
\usepackage{subcaption} % For sub-captions
\usepackage{makecell} %  Create multi-lined tabular cells
%--------------------------------------
% documents-specific packages
\usepackage{tikz} % For phase diagrams
\usetikzlibrary{arrows.meta} % Arrows library
\usepackage{pgfplots} % For native figures
\pgfplotsset{compat=1.17}
%--------------------------------------
\usepackage[footnote]{acronym} % For removing footnote indents
\usepackage{scrextend} % For removing footnote indents
\deffootnote[1em]{0em}{0em} % For removing footnote indents
{\textsuperscript{\thefootnotemark}\,\,}
\usepackage{multicol}
\usepackage{listings}
%--------------------------------------
\makeatother
\DeclarePairedDelimiter\ceil{\lceil}{\rceil}
\DeclarePairedDelimiter\floor{\lfloor}{\rfloor}
\DeclareMathOperator*{\argmax}{arg \ max}
%--------------------------------------
\renewcommand{\floatpagefraction}{.9}
\renewcommand{\topfraction}{.9}
\renewcommand{\bottomfraction}{.95}
\renewcommand{\textfraction}{.05}
%--------------------------------------
\newenvironment{myfont}{\fontfamily{qcr}\selectfont}{\par}
%--------------------------------------
% Some custom theorems
\newtheorem{theorem}{Theorem}
\newtheorem{prop}{Proposition}
\newtheorem{axiom}[theorem]{Axiom}
\newtheorem{case}[theorem]{Case}
\newtheorem{claim}[theorem]{Claim}
\newtheorem{conclusion}[theorem]{Conclusion}
\newtheorem{condition}[theorem]{Condition}
\newtheorem{conjecture}[theorem]{Conjecture}
\newtheorem{corollary}{Corollary}
\newtheorem{criterion}[theorem]{Criterion}
\newtheorem*{definition}{Definition}
\newtheorem{example}{Example}
\newtheorem{exercise}[theorem]{Exercise}
\newtheorem{lemma}{Lemma}
\newtheorem{notation}[theorem]{Notation}
\newtheorem{problem}[theorem]{Problem}
\newtheorem{proposition}[theorem]{Proposition}
\newtheorem{remark}{Remark}
\newtheorem{solution}[theorem]{Solution}
\newtheorem*{proprestate}{Proposition~1}
\newtheorem{summary}[theorem]{Summary}
%--------------------------------------
\newtheorem{innercustomthm}{Proposition}
\newenvironment{customthm}[1]
{\renewcommand\theinnercustomthm{#1}\innercustomthm}
{\endinnercustomthm}
\newtheorem{innercustomcor}{Corollary}
\newenvironment{customcor}[1]
{\renewcommand\theinnercustomcor{#1}\innercustomcor}
{\endinnercustomcor}
\newtheoremstyle{named}{}{}{\itshape}{}{\bfseries}{.}{.5em}{\thmnote{#3}#1}
\theoremstyle{named}
\newtheorem*{namedtheorem}{}
%--------------------------------------
\newcommand\independent{\protect\mathpalette{\protect\independenT}{\perp}}
\def\independenT#1#2{\mathrel{\rlap{$#1#2$}\mkern2mu{#1#2}}} % Define statistical independence symbol
%--------------------------------------
\usepackage[indent=0pt]{parskip} % for paragraph spacing
%--------------------------------------
\usepackage{comment} % for commenting long sections
%--------------------------------------
\raggedbottom % for ragged bottom to paragraphs
%----------------------------------------------------------------
\onehalfspacing  % for 1.5x line height
\frenchspacing

%----------------------------------------------------------------
\title[A model of optimal science funding]{A model of optimal science funding:\\ targeting the capability-resource gap}

\author{Aydin Mohseni}
\address{Department of Philosophy, Carnegie Mellon University}
\email{amohseni@andrew.cmu.edu}
\thanks{This work was supported by the John Templeton Foundation [grant number TODO].}

\author{Simon DeDeo}
\address{Department of Social and Decision Sciences, Carnegie Mellon University}

\author{Kevin J.\,S. Zollman}
\address{Department of Philosophy, Carnegie Mellon University}

\begin{document}

\begin{abstract}
The grant funder's challenge is to maximize the impact of a fixed budget under uncertainty about the researchers who might receive it. We develop a Bayesian model of this challenge: researchers differ in their capabilities and in their resources; research output requires both, and funding compounds capability across rounds. We show that the output-maximizing allocation of a round's budget funds the capability-resource gap: grants increase with capability and decrease with existing resources. Track records underdetermine the gap, because output confounds the two; what a funder needs is a signal of capability separate from resources, and peer review is one source of such a signal. We find that the value of that signal increases with how unequally the gaps are distributed, which capability inequality mainly drives, and with the signal's informativeness, and that under heavy-tailed capability distributions even a noisy signal captures most of that value. Uniform seed grants and lotteries, which forgo targeting, lose little research output exactly where review is worth little, and most of what a partial lottery loses comes from discarding the review ranking rather than from chance. Two features of a field, the funder's budget relative to researchers' resources and the inequality of the gaps, locate a funder on a map of these results.
\end{abstract}

\maketitle

\noindent\textbf{Significance.} Research funders divide fixed budgets among researchers whose abilities and needs they cannot observe. Whether to concentrate funds or spread them, whether peer review is worth its cost, and whether lotteries should replace it are unsettled. We model funding as a sequential decision under uncertainty and show that the output-maximizing allocation funds the gap between each researcher's capability and their resources, that track records alone underdetermine that gap, and that what peer review is worth and how lotteries and seed grants perform turn on two features of a field: the funder's budget relative to researchers' resources, and how unequal the gaps are, which mostly reflects how unequal capability is. We show where representative funders fall on these two features.

The United States spent close to a trillion dollars on research and development in 2024, roughly a fifth of it funded by the federal government \citep{NCSES2026, NSB2026}, and how best to allocate that money remains unsettled. Whether funding should be concentrated among a few researchers or spread across many is a long-running dispute \citep{Aagaard2020}. How well peer review predicts research outcomes is contested: among funded NIH grants, percentile scores predicted subsequent productivity barely better than chance \citep{Fang2016}, while studies with other designs find modest but real predictive power \citep{LiAgha2015, ParkLeeKim2015}. Several funders now allocate part of their budgets by lottery, in forms that range from a tie-break among top-ranked proposals to a draw among all applicants above a threshold \citep{Liu2020, Heyard2022, Feliciani2024}. And a randomized trial of winning a grant found no clear effect on the recipients' publications over the following four years \citep{Barnett2024}, which raises the prior question of when a marginal grant adds output at all.

These disputes concern one problem: what a funder with a fixed budget and imperfect knowledge of researchers should do. The empirical studies each measure one input to that decision, how well scores predict outcomes among funded grants or how research output responds to grant size, and none can say what the funder should do with what it knows; that depends on what the funder is uncertain about, what its information is worth, and how the answer changes with the budget. We treat the funder's problem as a sequential decision under uncertainty and analyze it using a Bayesian approach. Researchers differ in capability, what they know and can do, and in resources, what they have to work with; output requires both and is limited by the scarcer of the two; the funder observes output, and the noisy judgments of reviewers, but neither attribute directly. A grant adds resources in the round it is given and, because doing research builds expertise, raises capability in later rounds. Each round, the funder observes the output that results, updates its beliefs about the researchers, and allocates again.

Three results follow. First, the allocation of a round's budget that maximizes expected research output funds the gap between each researcher's capability and their resources: grants increase with capability and decrease with existing resources, and the budget sets how many researchers are funded at all. Two heuristics a funder might use, funding the track record and funding the under-resourced, each track one of the two quantities, and each can perform worse than simply dividing the budget uniformly across all researchers. Second, track records underdetermine the gap, because output confounds capability with resources; what the funder needs is a signal of capability separate from resources, and peer review is one source of such a signal. The value of that signal, and how seed grants and lotteries perform relative to targeted funding, depend on the same two features of a field: how much of researchers' resources the funder's budget can supply, and how unequally the capability-resource gaps are distributed, which in every setting we ran is mostly a matter of how unequal capability is. Where the budget is tight and capability heavy-tailed, targeting is worth a great deal, even noisy review captures most of that value, and lotteries give up much of it; where the budget is ample and capability evenly spread, little research output is at stake and spreading funds widely loses almost nothing. Within the model, no funding mechanism is best in every field. Third, the two features can be estimated from what funders already know, and we place representative funding programs on the map they define.

Two assumptions bound these results: grants are used up in the round they are given, with whatever lasts from them entering as capability, and capability and resources are complements. The objective is aggregate expected research output; diversity of directions, risk, and fairness are outside it, and ``optimal'' throughout means output-maximizing under that objective. The model is deliberately simple. Its contribution is to the theoretical understanding of the funder's problem: it shows that the disputes turn on two features of a field, and it says which way each dispute goes in each kind of field.

\section*{Results}

\subsection*{The optimal allocation funds the capability-resource gap}
Consider first a funder that knows every researcher's capability $K_i$ and resources $R_i$ and divides one round's budget $B$ so as to maximize that round's expected output, which for researcher $i$ is $AK_iR_i/(K_i + R_i)$, where $A$ is a productivity constant that sets the units of output and $R_i$ is the sum of baseline resources $R_{i0}$ and the grant $g_i$ (Materials and Methods). The marginal value of a dollar to researcher $i$ is $AK_i^2/(K_i + R_{i0} + g_i)^2$: it increases with their capability and decreases with the resources they already hold, grants included. An allocation is optimal exactly when no dollar can be moved to a researcher who would produce more with it, and solving that condition gives the gap rule (\ref{app:gap-rule}, Proposition~1):
\[
g_i^* = \max(c\,K_i - R_{i0},\, 0),
\]
where $c > 0$ is the value at which the grants sum to the budget. Each researcher's target is $c$ times their capability; the rule fills each positive gap between target and resources and grants nothing to the rest. The rule does not depend on the harmonic form. For any production function with constant returns to scale and diminishing returns to resources, $F(K, R) = A\,K\,h(R/K)$ with $h$ increasing and strictly concave, equal marginal value among the funded means a common ratio of resources to capability, and the same rule follows with $c$ set by the budget (\ref{app:gap-rule}, Proposition~1); the harmonic mean is the instance we simulate. The condition $R_{i0} < cK_i$ draws a frontier through the population, separating the funded from the unfunded (Fig.~\ref{fig:p1}). When the budget is small, $c$ is small and few researchers are funded; a larger budget raises $c$, moves the frontier outward, and funds more researchers with larger grants, until at a large enough budget funding extends across nearly the whole population.

The gap rule explains why each heuristic fails. Funding the track record fails because output reflects resources as well as capability: a productive, well-resourced researcher has a small or negative gap, and an additional dollar adds little to their output. Funding the under-resourced fails for the opposite reason: scarcity is no evidence of capability, and dollars given to researchers of low capability add little output however scarce their resources. Each heuristic can produce less output than dividing the budget equally (\ref{app:gap-rule}, Examples 1 and 2; Fig.~\ref{fig:heuristics} shows both schemes on the population of Fig.~\ref{fig:p1}).

How much targeting matters depends on the budget. Expected output increases with resources at a diminishing rate and never exceeds $AK_i$, so when the budget is large enough to bring every researcher's output near that bound, the optimal allocation and equal division differ little. The value of targeting, the expected output the optimal allocation adds over equal division, vanishes as the budget grows (\ref{app:gap-rule}, Corollary~\ref{cor:targeting-vanishes}), and in every one of several hundred simulated populations it also fell with the budget throughout, measured as a share of uniform funding's own gain (Fig.~\ref{fig:targeting-value}). Targeting matters where the budget is small relative to researchers' resources, and this dependence on the budget recurs in every result below.

Proposition~1 concerns one round. Over several rounds a grant also raises the recipient's future capability, and the allocation that maximizes output over the whole horizon need not take the gap form. Numerically it nearly does. Over two rounds with complete information, the exact dynamic optimum, free to choose its recipients and how much of the budget to spend in each round, exceeds the gap rule applied round by round with an even split by at most 1.4\% of the research output that funding adds, across every capability tail, compounding rate, and budget we sweep, and its first-round grants correlate with the gap rule's at 0.98 or above (\ref{app:gap-rule}, Table~\ref{tab:dynopt}). We therefore use the round-by-round gap rule as the complete-information benchmark throughout, and "optimal allocation" below refers to it.

\begin{figure}[p]
\centering
\begin{tikzpicture}
\begin{axis}[
  width=0.82\textwidth, height=7.2cm,
  axis lines*=left,
  xmin=0, xmax=8.6, ymin=0, ymax=6.5,
  xtick=\empty, ytick=\empty,
  x axis line style={-{Stealth[length=4pt]}, line width=0.4pt},
  y axis line style={-{Stealth[length=4pt]}, line width=0.4pt},
  xlabel={resources $R$},
  ylabel={capability $K$},
  ylabel style={font=\small}, xlabel style={font=\small},
  tick label style={font=\small},
  axis line style={line width=0.4pt},
  tick style={line width=0.4pt, black},
  clip=false,
]
% funding frontier K = R/c
\addplot[black, line width=0.5pt, domain=0:6.4, samples=2] ({x*0.8332}, {x});
\node[font=\small, anchor=west, align=left] at (axis cs:5.6,6.25)
  {funding frontier ($R = cK$)};
% grants: horizontal arrows from (R0, K) to (cK, K)
\addplot[quiver={u=\thisrow{u}, v=\thisrow{v}}, -{Stealth[length=3.5pt]},
         gray, line width=0.5pt] table {fig1-funded.dat};
% funded researchers (filled)
\addplot[only marks, mark=*, mark size=1.6pt, black] table[x=R, y=K] {fig1-funded.dat};
% unfunded researchers (open)
\addplot[only marks, mark=o, mark size=1.6pt, gray, line width=0.5pt]
  table[x=R, y=K] {fig1-unfunded.dat};
% direct labels
\node[font=\small, anchor=west] at (axis cs:0.4,5.2) {funded};
\node[font=\small, anchor=west] at (axis cs:6.0,2.6) {unfunded};
\end{axis}
\end{tikzpicture}
\caption{The optimal allocation funds the capability-resource gap. Each point is a researcher, placed by resources and capability; the line is the funding frontier $R = cK$ at budget scale $b = 0.1$. Researchers below the frontier (filled) are funded, and each arrow is a grant, which moves its researcher exactly to the frontier. Researchers on or beyond the frontier (open) receive nothing.}
\label{fig:p1}
\end{figure}


\subsection*{What the funder can learn: track records and peer review}
A real funder observes research output, not capability or resources. Expected output is one number produced by the two jointly, so a researcher of high capability and modest resources and a researcher of modest capability and abundant resources can produce identical records, and these are exactly the researchers the gap rule treats most differently: the first may have the largest gap in the population, the second none. Even expected output gives only a lower bound on capability: since expected output never exceeds $AK_i$, a researcher whose expected output is $\lambda_i$ has $K_i > \lambda_i/A$, and realized output, a Poisson draw around $\lambda_i$, bounds it less. The problem is worst where the gap is largest. When resources are small relative to capability, expected output is approximately $AR_i$, proportional to resources and nearly independent of capability (Fig.~\ref{fig:p2}\textit{A}), so on small grants output carries little information about capability, and only grants comparable to capability itself make records informative about it.

A Bayesian funder relying on records alone therefore allocates far from the gap rule, and further rounds of records improve its allocation only slowly: its expected output falls short of the complete-information optimum by 46\% of what that optimum adds in the first round and still by 19\% in round twenty, while a funder that also holds a single review score for each researcher, observed once at the start, is 8\% short in round one and 3\% by round twenty (Fig.~\ref{fig:p2}\textit{B}); in the same runs, the correlation of the records-only funder's grants with the optimal grants rises from 0.08 to 0.32, against 0.83 for the review funder in round one. The reason lies in the model's structure. Capabilities compound while resources do not, so each researcher's output approaches the level their resources sustain, and output at that level carries less and less information about capability.

What the funder lacks is a signal of capability separate from resources. Peer review can provide one: a noisy signal, but one whose noise does not come from resources. We model it as each researcher's capability plus Gaussian noise, observed once (Materials and Methods). We measure the value of review as the additional expected output of a funder allocating with the signal over the same funder allocating without it, and we report it as the share of the records-only shortfall that the signal recovers, where the records-only shortfall is the difference between the records-only funder's output and the output a funder with complete information would produce. That value depends on the field (Fig.~\ref{fig:p2}\textit{C}). Review is worth most where capability is heavy-tailed and the budget is tight, and under heavy-tailed capability even a signal of moderate informativeness captures most of its value: at the default signal noise, review retains 82\% of its maximum value when the capability distribution's tail parameter is $\alpha_K = 1.3$, 72\% at the default distribution ($\alpha_K = 2$), and 45\% where capability is spread more evenly ($\alpha_K = 3.5$); at three times the default noise, 52\%, 33\%, and 8\%. The three distributions have the same mean capability, so the tail parameter varies how dispersed and how heavy-tailed capability is, not its level. The reason is twofold. Where capability is heavy-tailed, most of the output a funder can add comes from getting resources to the few researchers of unusually high capability, and precisely those researchers are the easiest to identify: they stand far apart from everyone else, so even a noisy signal separates them. Where capability is spread more evenly, no researcher matters much more than another, so review of any informativeness adds little; and where the budget is ample, information about whom to fund has little left to add. This ordering of review's value by field holds across production functions from Cobb-Douglas to Leontief (\ref{app:technology}), and it survives the checks a skeptic would ask for: it holds when review is compared across fields at equal rank correlation between score and capability rather than at equal noise, for applicant pools of 25 to 200, for heavier and lighter tails of the resource distribution, when the review score partly reflects resources rather than capability alone, when review noise is multiplicative rather than additive, when review is repeated afresh every round rather than observed once, and across resource signals from nearly exact to nearly uninformative (\ref{si:sensitivity}). Of these variants, multiplicative review noise, a constant relative error, is the one that narrows the heavy-tailed field's advantage, because a constant relative error is a large absolute error for the researchers who matter most there: at equal informativeness review recovers 0.66 of the shortfall in the heavy-tailed field against 0.49 in the evenly spread field, where the additive signal gives 0.87 against 0.55. What review supplies, in every one of these variants, is information about capability that output cannot supply on its own; peer review is one source of such information, and the results apply to any other, such as a pilot grant large enough to make output informative. One feature changes how much any of this is worth: the association between capability and resources. Where the capable already hold the resources, the gaps are small and even, and both targeting and review are worth less; where resources sit with the less capable, both are worth more. One condition qualifies it: the funder must weight review at its true informativeness. A funder that overtrusts review, treating a noisy signal as reliable, can do worse than one that ignores review altogether, and overtrusting review loses more output than undertrusting it (Fig.~\ref{fig:overtrust}).

\begin{figure}[tp]
\centering
\begin{minipage}{\textwidth}\noindent\textbf{A}\par\vspace{-1.2em}
\begin{tikzpicture}
\begin{axis}[
  width=0.82\textwidth, height=6.0cm,
  axis lines*=left,
  xmin=0, xmax=6.6, ymin=0, ymax=10.8,
  xtick=\empty, ytick=\empty,
  x axis line style={-{Stealth[length=4pt]}, line width=0.4pt},
  y axis line style={-{Stealth[length=4pt]}, line width=0.4pt},
  xlabel={resources $R$},
  ylabel={expected output},
  ylabel style={font=\small}, xlabel style={font=\small},
  axis line style={line width=0.4pt},
  clip=false,
]
\addplot[black, line width=0.7pt] table[x=R, y=k30] {fig3-curves.dat};
\addplot[black, line width=0.7pt] table[x=R, y=k10] {fig3-curves.dat};
\addplot[black, line width=0.7pt] table[x=R, y=k3]  {fig3-curves.dat};
\node[font=\small, anchor=west] at (axis cs:6.7,10.0) {$K = 30$};
\node[font=\small, anchor=west] at (axis cs:6.7,7.5)  {$K = 10$};
\node[font=\small, anchor=west] at (axis cs:6.7,4.0)  {$K = 3$};
\end{axis}
\end{tikzpicture}
\end{minipage}\par\vspace{0.9em}
\begin{minipage}{\textwidth}\noindent\textbf{B}\par\vspace{-1.2em}
\begin{tikzpicture}
\begin{axis}[
  width=0.82\textwidth, height=6.0cm,
  axis lines*=left,
  xmin=1, xmax=20, ymin=0, ymax=1,
  xtick={1, 5, 10, 15, 20},
  ytick={0, 0.25, 0.5, 0.75, 1},
  yticklabels={0, 0.25, 0.5, 0.75, 1},
  xlabel={round},
  ylabel={correlation with optimal grants},
  ylabel style={font=\small}, xlabel style={font=\small},
  tick label style={font=\small},
  axis line style={line width=0.4pt},
  tick style={line width=0.4pt, black},
  clip=false,
]
\addplot[black, line width=0.7pt] table[x=round, y=s5] {fig5-rounds.dat};
\addplot[black, line width=0.7pt] table[x=round, y=s4] {fig5-rounds.dat};
\node[font=\small, anchor=west] at (axis cs:20.4,0.90) {with review};
\node[font=\small, anchor=west] at (axis cs:20.4,0.41) {records only};
\end{axis}
\end{tikzpicture}
\end{minipage}\par\vspace{0.9em}
\begin{minipage}{\textwidth}\noindent\textbf{C}\par\vspace{-1.2em}
\begin{tikzpicture}
\begin{axis}[
  width=0.82\textwidth, height=6.0cm,
  axis lines*=left,
  xmode=log,
  xmin=0.05, xmax=20, ymin=0, ymax=1,
  xtick={0.05, 0.3, 1, 3, 10, 20},
  xticklabels={0.05, 0.3, 1, 3, 10, 20},
  ytick={0, 0.25, 0.5, 0.75, 1},
  yticklabels={0, 0.25, 0.5, 0.75, 1},
  xlabel={signal noise $\tau_K$},
  ylabel={share of shortfall recovered},
  ylabel style={font=\small}, xlabel style={font=\small},
  tick label style={font=\small},
  axis line style={line width=0.4pt},
  tick style={line width=0.4pt, black},
  clip=false,
]
\addplot[black, line width=0.7pt] table[x=tau, y=s13] {fig4-value.dat};
\addplot[black, line width=0.7pt] table[x=tau, y=s20] {fig4-value.dat};
\addplot[black, line width=0.7pt] table[x=tau, y=s35] {fig4-value.dat};
\node[font=\small, anchor=west] at (axis cs:22,0.31)  {$\alpha_K = 1.3$};
\node[font=\small, anchor=west] at (axis cs:22,0.115) {$\alpha_K = 2$};
\node[font=\small, anchor=west] at (axis cs:22,0.005) {$\alpha_K = 3.5$};
\end{axis}
\end{tikzpicture}
\end{minipage}\par\vspace{0.9em}
\caption{What the funder can learn. (\textit{A}) Small grants are uninformative about capability: expected output against resources for researchers of capability $K = 3$, $10$, and $30$; the curves coincide where resources are small relative to capability. (\textit{B}) Records do not catch up: the correlation of each funder's grants with the optimal grants, round by round; the records-only funder improves slowly and after twenty rounds remains further from the optimal allocation than the funder with the review signal is after one. (\textit{C}) Review is worth most where capability is heavy-tailed: the share of the records-only shortfall that the review signal recovers, as review noise $\tau_K$ increases, for three capability distributions. Settings: SI Text 7.}
\label{fig:p2}
\end{figure}


\subsection*{Spreading funds: seed grants, lotteries, and concentration}
Two alternatives to targeting recur in funding policy, and each has a rationale. Uniform seed grants give every researcher the same amount, so that a researcher whom review misjudges still receives something; in practice a funder gives out part of its budget as seed grants and targets the rest. Lotteries give grants to applicants chosen at random, usually from among those who pass a review threshold; their advocates argue that review cannot reliably rank the applicants it deems fundable, that review is costly, and that review is biased \citep{FangCasadevall2016, Roumbanis2019}. Both alternatives give up part of targeting on purpose, so both give up part of the value of targeting, and that value depends on the field. For a funder allocating with records and review, targeting adds more research output than uniform funding's entire gain over no funding in a heavy-tailed field with reliable review and a tight budget, and about a quarter of that gain in the default field at the largest budget we sweep (\ref{si:extended}). What the alternatives save, in the cost of review and in the effort researchers spend on proposals, lies outside our model \citep{GrossBergstrom2019}.

Consider seed grants first. The output lost increases faster than the fraction of the budget given out as seed grants: giving out half of the budget as seed grants loses more than twice the output that giving out a quarter loses (Fig.~\ref{fig:p4}\textit{A}). The size of the loss is set by what targeting is worth in the field: when a review-informed funder gives out three quarters of its budget as seed grants, it forfeits about 11\% of the research output its funding adds in a heavy-tailed field with reliable review (17\% at the smallest budget we sweep), and about 3.5\% in the default field.

Now consider lotteries. The lotteries funders run are partial: review sets a fundable pool, and chance picks the grantees within it. Such a lottery gives up two things, the review ranking within the pool and the certainty of who is funded, and the question is what each costs. We evaluate five schemes in a single round (Fig.~\ref{fig:p4}\textit{B}): equal division among the fifth of researchers with the highest review scores (selection by review); equal division among the top two fifths; the partial lottery that some funders run, which funds half of the top two fifths at random with full-size grants; uniform funding; and a lottery over all researchers with full-size grants. In a heavy-tailed field with a tight budget, selection by review gains about three quarters of what the optimal allocation gains when review is reliable and about half when review is very noisy. The partial lottery gains less, and most of the difference comes from ignoring the review ranking within the top two fifths rather than from the randomness itself: at the default noise, the partial lottery falls 0.18 short of selection by review, of which ignoring the ranking accounts for 0.14 and the randomness for 0.04. The lottery over all researchers gains least, less than uniform funding. In an evenly spread field with a large budget, the ordering reverses: uniform funding gains more than four fifths of what the optimal allocation gains, and every scheme that concentrates grants, whether by review score or by chance, gains less. A lottery among applicants above a review threshold therefore loses little output only where the review ranking it ignores is worth little: replacing that ranking by chance loses a few hundredths of the optimal allocation's gain in an evenly spread field or when review is very noisy, and about a sixth of that gain in a heavy-tailed field with reliable review. In an evenly spread field, the better alternative to review is small grants to many researchers; a lottery among a few applicants above a review threshold produces less.

The randomness itself costs output for a general reason. A lottery gives each applicant in its pool a chance at a full grant instead of a certain smaller one, and on average the applicant receives the same amount either way; because expected output increases with resources at a diminishing rate, the certain grant produces more. Dividing a pool's money equally among the pool therefore produces more expected output than any lottery over the same pool, provided that grants are divisible and that output rises with resources at a diminishing rate at every grant size (\ref{app:lottery}, Proposition~\ref{prop:lottery}). Where a project needs a minimum viable grant, the first condition fails and a lottery over full awards can beat dividing the same money into grants too small to use (\ref{app:lottery}). The cost of the randomness is small in our simulations; the cost of the discarded ranking is the one that depends on the field.

Seed grants and lotteries both spread funds more widely than targeting does, and both raise the older question of whether funding should be concentrated among a few researchers or spread across many. The production function does not settle that question. Stronger complementarity between capability and resources concentrates the review-informed funder's grants when the budget is tight in a heavy-tailed field, spreads them when the budget is ample, and has a non-monotone effect where capability is evenly spread (Fig.~\ref{fig:concentration}; \ref{si:extended}). How widely funds should be spread is set jointly by the production function, the budget, and the field's capability distribution.

\begin{figure}[tp]
\centering
\begin{minipage}{\textwidth}\noindent\textbf{A}\par\vspace{-0.3em}
\centering
\begin{tikzpicture}
\begin{axis}[
  width=0.76\textwidth, height=5.4cm,
  axis lines*=left,
  xmin=0, xmax=0.75, ymin=0, ymax=20,
  xtick={0, 0.25, 0.5, 0.75},
  ytick={0, 5, 10, 15, 20},
  xlabel={fraction of the budget given out as seed grants},
  ylabel={research output lost (\% of the gain from funding)},
  ylabel style={font=\small}, xlabel style={font=\small},
  tick label style={font=\small},
  axis line style={line width=0.4pt},
  tick style={line width=0.4pt, black},
  clip=false,
]
\addplot[black, line width=0.7pt, mark=*, mark size=1.2pt] table[x=xseed, y=hb1]  {fig10-floor-cost.dat};
\addplot[black, line width=0.7pt, mark=*, mark size=1.2pt] table[x=xseed, y=hb05] {fig10-floor-cost.dat};
\addplot[black, line width=0.7pt, mark=*, mark size=1.2pt] table[x=xseed, y=hb01] {fig10-floor-cost.dat};
\addplot[black, line width=0.7pt, mark=*, mark size=1.2pt] table[x=xseed, y=db05] {fig10-floor-cost.dat};
\node[font=\small, anchor=west] at (axis cs:0.77,8.11) {heavy-tailed, $b = 2$};
\node[font=\small, anchor=west] at (axis cs:0.77,11.07) {heavy-tailed, $b = 1$};
\node[font=\small, anchor=west] at (axis cs:0.77,17.29) {heavy-tailed, $b = 0.2$};
\node[font=\small, anchor=west] at (axis cs:0.77,3.47) {intermediate, $b = 1$};
\end{axis}
\end{tikzpicture}
\end{minipage}\par\vspace{0.9em}
\begin{minipage}{\textwidth}\noindent\textbf{B}\par\vspace{-0.3em}
\centering
\begin{tikzpicture}
\begin{axis}[
  name=left,
  width=0.36\textwidth, height=5.6cm,
  axis lines*=left,
  xmode=log,
  xmin=0.3, xmax=10, ymin=0, ymax=1,
  xtick={0.3, 1, 3, 10},
  xticklabels={0.3, 1, 3, 10},
  ytick={0, 0.25, 0.5, 0.75, 1},
  yticklabels={0, 0.25, 0.5, 0.75, 1},
  xlabel={review noise $\tau$},
  ylabel={share of the optimal allocation's gain},
  ylabel style={font=\small}, xlabel style={font=\small},
  tick label style={font=\small},
  axis line style={line width=0.4pt},
  tick style={line width=0.4pt, black},
  title={\small heavy-tailed capabilities,\\ tight budget},
  title style={at={(0,1)}, anchor=south west, xshift=-2pt, yshift=-1pt, align=left},
  clip=false,
]
\addplot[black, line width=0.7pt, mark=*, mark size=1.1pt] table[x=tau, y=ss] {fig11-heavy.dat};
\addplot[black, line width=0.7pt, mark=o, mark size=1.3pt] table[x=tau, y=ws] {fig11-heavy.dat};
\addplot[black, line width=0.7pt, mark=triangle*, mark size=1.4pt] table[x=tau, y=sl] {fig11-heavy.dat};
\addplot[black!65, dashed, line width=0.6pt, domain=0.3:10] {0.445};
\addplot[black!65, dotted, line width=0.8pt, domain=0.3:10] {0.378};
\end{axis}
\begin{axis}[
  name=right,
  at={(left.east)}, anchor=west, xshift=30pt,
  width=0.36\textwidth, height=5.6cm,
  axis lines*=left,
  xmode=log,
  xmin=0.3, xmax=10, ymin=0, ymax=1,
  xtick={0.3, 1, 3, 10},
  xticklabels={0.3, 1, 3, 10},
  ytick={0, 0.25, 0.5, 0.75, 1},
  yticklabels={,,,,},
  xlabel={review noise $\tau$},
  xlabel style={font=\small},
  tick label style={font=\small},
  axis line style={line width=0.4pt},
  tick style={line width=0.4pt, black},
  title={\small evenly spread capabilities,\\ ample budget},
  title style={at={(0,1)}, anchor=south west, xshift=-2pt, yshift=-1pt, align=left},
  clip=false,
  legend style={draw=none, fill=none, font=\scriptsize, at={(1.04,0.55)}, anchor=west, legend cell align=left, row sep=1pt},
  legend image post style={mark size=1.1pt},
]
\addplot[black, line width=0.7pt, mark=*, mark size=1.1pt] table[x=tau, y=ss] {fig11-even.dat};
\addlegendentry{top fifth by review}
\addplot[black, line width=0.7pt, mark=o, mark size=1.3pt] table[x=tau, y=ws] {fig11-even.dat};
\addlegendentry{top two fifths by review}
\addplot[black, line width=0.7pt, mark=triangle*, mark size=1.4pt] table[x=tau, y=sl] {fig11-even.dat};
\addlegendentry{partial lottery}
\addplot[black!65, dashed, line width=0.6pt, domain=0.3:10] {0.831};
\addlegendentry{uniform funding}
\addplot[black!65, dotted, line width=0.8pt, domain=0.3:10] {0.414};
\addlegendentry{lottery over all}
\end{axis}
\end{tikzpicture}
\end{minipage}\par\vspace{0.9em}
\caption{Seed grants and lotteries. (\textit{A}) The research output lost when a review-informed funder gives out a fraction of its budget as equal grants to all researchers and targets the rest, as a percent of the research output its grants add over no funding, in four fields and budgets. (\textit{B}) Each of five funding schemes' gain over no funding as a share of the optimal allocation's gain, as review noise grows, in a heavy-tailed field with a tight budget (left) and an evenly spread field with an ample budget (right). Selection by review divides the budget equally among the top fifth or the top two fifths by review score; the partial lottery gives full-size grants to a random half of the top two fifths. Settings: \ref{app:simulation}.}
\label{fig:p4}
\end{figure}


\section*{Discussion}
Our results isolate one quantity as key to optimal grant funding: the value of targeting researchers with the greatest gap between their capabilities and their resources. Two features of a field determine this value. The first is the funder's budget relative to researchers' own resources: where the budget is small, dollars are scarce and where they go matters; where the budget can relieve most researchers' bottlenecks, the optimal allocation and equal division produce nearly the same research output. The second is how unequally the capability-resource gaps are distributed. Across every setting we ran, the value of targeting is close to a function of the inequality of the optimal grants alone (\ref{si:sensitivity}), and what makes the gaps unequal is mainly how unequal capability is: where a few researchers can produce far more than the rest, getting resources to those few is most of what a funder can add. The association between capability and resources works through the same channel: where the capable already hold the resources, the gaps are small and even, and targeting is worth less. Every result above is a statement about this quantity: what records and review contribute to capturing it, and what seed grants and lotteries give up of it. Two results lie outside it. How widely a review-informed funder should spread its grants depends, in addition, on how strongly capability and resources complement each other (\ref{si:extended}, Fig.~\ref{fig:concentration}); and the timing of spending over the horizon, which we analyze in \ref{si:timing}, adds little research output beside the choice of recipients in every setting we test, subject to the limits stated there.

These results have a practical consequence for funders. Within the model, no funding mechanism is best in every field. Peer review, seed grants, and lotteries each perform better in some fields than in others, and the two features of the field decide which, so a funder choosing among mechanisms should begin by locating its field on the two (Fig.~\ref{fig:regime}). The first feature is observable. A funder's budget relative to the field's own resources can be read from administrative data: its annual awards to the field against the field's research spending from every other source. The second feature is not. The inequality of the gaps depends on capabilities the funder cannot see and on the budget itself, and the inequality of the optimal grants, the quantity that organizes our results (\ref{si:sensitivity}), is not one a funder can compute. Three observable quantities stand in for it. The inequality of research output in the field is the natural proxy for the inequality of capability, though where grants are small output tracks resources rather than capability, so this proxy is imperfect in the direction of the resource distribution. The dispersion of review scores over the whole applicant pool is a second, with the same caveat about noise. The most direct is the response of output to grants large enough to relieve the bottleneck: on a grant comparable to the recipient's capability, output becomes informative about capability (Fig.~\ref{fig:p2}\textit{A}), so a program of substantial pilot grants measures the gaps, whereas small seed grants do not, because on small grants output tracks resources for every recipient. How existing funding is distributed completes the estimate: where resources already sit with the capable, the gaps are smaller and more even than capability inequality alone would imply.

Fig.~\ref{fig:regime} places four kinds of funding program on the budget axis by order of magnitude, from 2023 and 2024 budget figures (Materials and Methods). A national biomedical agency that supplies a large share of its field's research support sits near a budget equal to the field's other resources; a national science agency that funds its fields more thinly sits near a tenth of that; a continental excellence council near a few hundredths; and a private foundation funding one field among many at a hundredth or below. Where a field lies on the inequality axis is the estimate each funder must make, and the range of tail parameters estimated for research productivity spans most of the axis, so the placements say which regime a funder's budget puts it in, not where it sits within the regime. The regimes differ sharply. A foundation or council with a small budget in a heavy-tailed field faces the largest value of targeting per dollar: review directed at identifying the few researchers of unusually high capability is worth most, even noisy review recovers most of that value, seed grants give up the most output, and a lottery among applicants above a review threshold gives up about a sixth of what selection by review gains. An agency whose budget approaches the field's other resources, in a field where capability is evenly spread, faces a small value of targeting: review of any informativeness adds little, seed grants and lotteries give up little, and the model's best use of the budget is small grants to many researchers.

% PNAS Fig. 4 (Discussion): the regime map, two panels; copy of fig16-regime.tex with SI refs. x = budget relative to the field's
% resources (b, log scale); y = capability inequality (Gini coefficient of the Pareto
% capability distribution, 1/(2 alpha_K - 1), with alpha_K in parentheses). (A) value of
% targeting: the complete-information optimum's gain over uniform funding as a share of
% uniform funding's gain. (B) value of review at the default noise, as a share of the research
% output that review-informed funding adds. Shaded bands: four kinds of program placed on the
% budget axis by order of magnitude (SI Materials and Methods). Data: fig16-targ.dat, fig16-rev.dat from
% regime_m.R (M = 1000) via make_regime.py; contour label positions computed from the data.
\begin{figure}[t]
\centering
\newcommand{\regimebands}{%
\fill[black!8] (axis cs:0.01,0.1) rectangle (axis cs:0.03,0.72);
\fill[black!8] (axis cs:0.04,0.1) rectangle (axis cs:0.08,0.72);
\fill[black!8] (axis cs:0.2,0.1) rectangle (axis cs:0.45,0.72);
\fill[black!8] (axis cs:1.5,0.1) rectangle (axis cs:3,0.72);
\node[font=\scriptsize, anchor=south, align=center] at (axis cs:0.0173,0.73) {foundation};
\node[font=\scriptsize, anchor=south, align=center] at (axis cs:0.0566,0.79) {excellence\\council};
\node[font=\scriptsize, anchor=south, align=center] at (axis cs:0.3,0.73) {science\\agency};
\node[font=\scriptsize, anchor=south, align=center] at (axis cs:2.12,0.73) {biomedical\\agency};}
\begin{tikzpicture}
\begin{axis}[
  name=panelA,
  scale only axis, width=4.5cm, height=5.4cm,
  axis lines*=left,
  xmode=log,
  xmin=0.01, xmax=3, ymin=0.1, ymax=0.72,
  xtick={0.01, 0.03, 0.1, 0.3, 1, 3},
  xticklabels={0.01, 0.03, 0.1, 0.3, 1, 3},
  ytick={0.111, 0.167, 0.25, 0.333, 0.5, 0.625, 0.714},
  xlabel={budget relative to the field's resources, $b$},
  ylabel style={font=\small}, xlabel style={font=\footnotesize, at={(rel axis cs:1.14,-0.1)}, anchor=north},
  tick label style={font=\small},
  axis line style={line width=0.4pt},
  tick style={line width=0.4pt, black},
  clip=false,
  yticklabels={0.11 (5), 0.17 (3.5), 0.25 (2.5), 0.33 (2), 0.50 (1.5), 0.63 (1.3), 0.71 (1.2)},
  ylabel={capability inequality: Gini coefficient ($\alpha_K$)},
]
\regimebands
\addplot[black, line width=0.6pt, empty line=jump] table[x expr=10^\thisrow{x}, y=y] {fig16-targ.dat};
\node[font=\scriptsize, anchor=west] at (axis cs:3,0.178) {0.1};
\node[font=\scriptsize, anchor=west] at (axis cs:3,0.304) {0.25};
\node[font=\scriptsize, anchor=west] at (axis cs:3,0.446) {0.5};
\node[font=\scriptsize, anchor=west] at (axis cs:3,0.639) {1};
\node[font=\scriptsize, fill=white, inner sep=1pt] at (axis cs:0.195,0.60) {2};
\node[font=\scriptsize, fill=white, inner sep=1pt] at (axis cs:0.032,0.675) {3};
\node[font=\bfseries, anchor=south west] at (rel axis cs:-0.28,1.13) {A};
\end{axis}
\begin{axis}[
  at={(panelA.east)}, anchor=west, xshift=1.3cm,
  scale only axis, width=4.5cm, height=5.4cm,
  axis lines*=left,
  xmode=log,
  xmin=0.01, xmax=3, ymin=0.1, ymax=0.72,
  xtick={0.01, 0.03, 0.1, 0.3, 1, 3},
  xticklabels={0.01, 0.03, 0.1, 0.3, 1, 3},
  ytick={0.111, 0.167, 0.25, 0.333, 0.5, 0.625, 0.714},
  xlabel={},
  ylabel style={font=\small}, xlabel style={font=\footnotesize},
  tick label style={font=\small},
  axis line style={line width=0.4pt},
  tick style={line width=0.4pt, black},
  clip=false,
  yticklabels={}, ylabel={},
]
\regimebands
\addplot[black, line width=0.6pt, empty line=jump] table[x expr=10^\thisrow{x}, y=y] {fig16-rev.dat};
\node[font=\scriptsize, anchor=west] at (axis cs:3,0.199) {0.05};
\node[font=\scriptsize, anchor=west] at (axis cs:3,0.277) {0.1};
\node[font=\scriptsize, anchor=west] at (axis cs:3,0.435) {0.2};
\node[font=\scriptsize, anchor=west] at (axis cs:3,0.693) {0.3};
\node[font=\scriptsize, fill=white, inner sep=1pt] at (axis cs:0.83,0.60) {0.4};
\node[font=\bfseries, anchor=south west] at (rel axis cs:-0.1,1.13) {B};
\end{axis}
\end{tikzpicture}
\caption{Where a funder stands. (\textit{A}) The value of targeting: the complete-information optimum's gain over uniform funding, as a share of uniform funding's gain over no funding. (\textit{B}) The value of review at the default noise, as a share of the research output that review-informed funding adds. Both by the funder's budget relative to the field's resources and the inequality of capability, at fixed mean capability; contours are labeled by value. Shaded bands place four kinds of program on the budget axis by order of magnitude, from 2023 and 2024 annual budget figures with $b$ equal to twice the annual ratio, since the map's budget covers two rounds (\ref{app:simulation}); where a field lies on the vertical axis is the estimate each funder must make, and the tail parameters estimated for research productivity lie in its upper part. Two rounds, $\tau_K = 1$, 50 populations per cell; contours interpolated between the 63 cells. Settings: \ref{app:simulation}.}
\label{fig:regime}
\end{figure}


Three further implications concern peer review. Since most of review's value comes from separating the few researchers of unusually high capability from the rest, and fine distinctions among the many applications of comparable merit add little research output, a review system that spends most of its effort ranking the middle of the distribution spends that effort where the model locates the least output. Second, the finding that percentile scores barely predict productivity comes from funded NIH grants with scores in the top fifth \citep{Fang2016}, and the studies that find modest predictive power also use funded grants \citep{LiAgha2015}; both measure prediction among applicants that review has already selected, within the range where the model expects review to add least, so a weak association there is consistent with review of considerable value over the whole applicant pool. Third, a funder that treats review as more reliable than it is loses more output than a funder that treats review as noisier than it is, so a funder uncertain about review's accuracy should err toward undertrust.

The model holds several things fixed, and its results hold within those limits. There is one funder, and researchers do not respond strategically, so competition among funders, proposal effort, risk choice, and the cost of review are outside the model; the contest models of grant competition cover that ground \citep{GrossBergstrom2019, GrossBergstrom2025}, and the mechanism-design treatment of grants by \citet{Carnehl2024}, in which the funder must also elicit information from applicants who bear the cost of applying, is the natural complement to a model in which the funder's information arrives for free. \citet{Sorzano2026} also treat funding as a decision under heavy-tailed returns and find more room for lotteries than we do; their heavy tail is in observed impact, whereas ours decomposes output into latent capability and resources, and that decomposition is what gives review its value here. Output is a single number per round, so projects do not differ in kind and exploring new directions has no value in the model \citep{AvinBJPS2019}; nor do a researcher's proposals differ in quality, so this is a model of investigator-level funding, and project-specific quality that review might observe is the natural extension. Grants are divisible, so a minimum viable grant size, which can favor a lottery over full awards, is outside the model (\ref{app:lottery}). The review signal is noisy but unbiased; bias and conservatism in review, which are among the strongest arguments for lotteries, are not modeled. Capability is a state, not a trait: it holds whatever a researcher has accumulated that raises output, expertise, methods, a trained group, and doing research adds to it; a funder that reads it as innate ability misreads the model. The asymmetry between the two inputs is an assumption with consequences: capability persists and grows, resources recur and do not persist, so a grant's lasting effect enters only through capability. Durable equipment, data, and infrastructure violate this asymmetry. The one-round results do not depend on it; the results about compounding and about the timing of spending (\ref{si:timing}) do. The compounding here is of capability through doing research, and differs from the cumulative advantage in which winning a grant raises the chance of winning the next \citep{BolVaanRijt2018}; the two would need to be separated in any empirical test. Inputs are complements, and the timing results depend on complementarity while the information results do not. The numbers we report are comparisons across the parameter values we sweep, and none is calibrated to a real funder; what applies to a real field is the qualitative features of the dynamics that organize the results.

The model makes four predictions that data on funded research can test. The effect of winning a grant on output should be larger for recipients whose other resources are small relative to their capability, so a randomized trial of funding, such as the one that found no clear average effect \citep{Barnett2024}, should find its effect concentrated among under-resourced recipients, and a null average is what the model predicts where most recipients are already well resourced. Measures that stand in for capability, such as prior record or review score, should predict research output better the larger a grant is relative to the recipient's other resources, because on small grants output is nearly independent of capability. How well review scores predict output, measured over the whole applicant pool rather than among funded grants, should be higher in fields whose output distributions have heavier tails. And where funders run lotteries beside review, the output lost to the lottery should be larger in heavier-tailed fields. Each of the remaining predictions can be tested with data that funders already hold: the first by comparing how well review scores predict output across grant sizes within one applicant pool; the second by comparing that predictive power across fields, scored over the whole applicant pool where the outcomes of funded and unfunded applicants can both be observed; the third by comparing the outputs of lottery-funded and review-funded grants where a funder runs both.

A funder with a fixed budget and imperfect knowledge of researchers faces one problem, and our model gives it one answer. The output-maximizing allocation funds the gap between each researcher's capability and their resources. Records alone underdetermine that gap, because output does not identify which input binds; a signal of capability separate from resources does, and peer review is one source of it. What that signal is worth, and how seed grants and lotteries perform relative to targeted funding, turn on the same two features of a field, and the map of those features has regimes in which the answers differ. A funder deciding among review, lotteries, and seed grants should therefore begin by asking which kind of field it is funding.

\section*{Materials and Methods}
\textbf{Model.} A funder faces $n$ researchers over $T$ funding rounds. Researcher $i$ has capability $K_i$ and resources $R_i$; initial values are drawn from Pareto distributions with tail parameters $\alpha_K$ and $\alpha_R$ (smaller values give more unequal fields; below 2 the variance is infinite), coupled through a Gaussian copula with parameter $\rho$ (a Pearson correlation between Pareto variables is undefined when their variances are infinite; $\rho = 0.5$ and $0.8$ correspond to rank correlations of 0.48 and 0.79). Capability is a state variable: it holds whatever a researcher has accumulated that raises output, and doing research adds to it. In each round, researcher $i$ produces research output drawn from a Poisson distribution with mean $\lambda_i = AK_iR_i/(K_i + R_i)$, proportional to the harmonic mean of the two inputs, where the productivity constant $A$ sets the units of output (the simulations set $A = 1$); output increases in both inputs and is limited by the scarcer. Between rounds capability compounds, $K_i \leftarrow K_i + \epsilon\,\lambda_i/A$, at rate $\epsilon$. Resources do not accumulate: $R_i = R_{i0} + g_i$, where the baseline $R_{i0}$ recurs every round and the grant $g_i$ is consumed in the round it is given. The unit of analysis is the researcher, not a researcher-project pair. The funder's total budget over the horizon is $b\,n\,\mathbb{E}[R_0]$, so at budget scale $b = 1$ it equals the community's expected baseline resources for one round; the budget does not depend on the horizon.

\textbf{Information and strategies.} The funder knows the distributions but observes no researcher's capability or resources. It observes each researcher's realized output, a Poisson draw with mean $\lambda_i$, as it accumulates; a resource signal $R_{i0} + \eta_i$, observed once, with $\eta_i$ Gaussian with standard deviation $\tau_R$; and a review signal $K_i + \xi_i$, observed once at the start by the strategies that use it, with $\xi_i$ Gaussian with standard deviation $\tau_K$ and independent across researchers; a reviewer's misjudgment of a researcher therefore persists, and a variant with a fresh review every round is in \ref{si:sensitivity}. Bayesian strategies hold a posterior over each researcher's capability and baseline resources, computed from these likelihoods by importance sampling with $M$ particles drawn from the prior ($M = 1{,}000$ for Fig.~\ref{fig:p2}\textit{B} and \textit{C}, Fig.~\ref{fig:regime}, and the sensitivity analysis, and $200$ elsewhere; the convergence of the results in $M$ is checked in SI Materials and Methods), update it after each round, and allocate each dollar to the researcher whose expected output, given the posterior, rises most with it. They differ in information (with or without the review signal) and in planning. The myopic funder maximizes the current round's expected output. The forward-looking funder is a certainty-equivalent receding-horizon planner: each round it chooses the current allocation together with a plan for the remaining budget so as to maximize expected output over the remaining horizon, treating future output as equal to its expectation, and it re-plans after observing the round's output; it does not value the information that its own grants will generate (\ref{si:timing}). The complete-information benchmark applies the gap rule to each round's share of the budget with the true capabilities and resources. Three baselines complete the set: no funding, against which every result is normalized; uniform funding; and track-record funding, which allocates in proportion to the previous round's output. Seed grants and lotteries modify these strategies as described in Results (\ref{sec:model}, Table~\ref{tab:strategies}).

\textbf{Parameters and scoring.} Tail parameters range from 1.3 to 3.5, a range that includes the Pareto tails near 1 to 2 estimated for research productivity \citep{Lotka1926, Nicholls1989}; since productivity is an imperfect guide to capability (Results), the range extends well beyond those estimates; the budget scale $b$ from 0.01 to 3, with larger budgets in the timing analyses (\ref{si:timing}); the compounding rate $\epsilon$ from near zero to 0.85; review noise $\tau_K$ from 0.05 to 20; the horizon from one round to ten. Default values are $T = 2$, $n = 50$, $\epsilon = 0.1$, $b = 1$, $\alpha_K = \alpha_R = 2$, $\tau_K = \tau_R = 1$, $\rho = 0$. Each strategy is scored by the sum of expected outputs along its run, which equals realized output averaged over repetitions of the same allocations, and strategies are compared on identical simulated populations with identical signal and output draws, so comparisons are paired. Results are reported as shares of a named comparison, most often the research output that funding adds: a strategy's expected output minus the expected output under no funding. Analytic results (the gap rule, its corollaries, and the lottery proposition) are proved in \ref{app:gap-rule} and \ref{app:lottery}; the settings behind each figure are in \ref{app:simulation}. Code and data are available in the project repository. % [check: archived commit at publication]

\textbf{Placing programs on the budget axis.} The budget scale $b$ of a real program is its annual awards to a field divided by the field's annual research spending from every other source, which stands in for the baseline resources of the field's researchers; the estimate counts money and not time, and treats each funder's field as the disciplines it mainly funds, so it is an order of magnitude. United States universities spent \$108.8 billion on research and development in fiscal year 2023, \$55.2 billion of it in the life sciences, and received \$33.1 billion from the Department of Health and Human Services, chiefly the National Institutes of Health, and \$6.7 billion from the National Science Foundation \citep{NCSES2025HERD}: a biomedical agency of that size has $b$ near 1 (\$33 billion against \$22 billion from other sources), and a science agency of that size funding the other half of academic research has $b$ near 0.15. The European Research Council's budget of 16 billion euros over 2021 to 2027 \citep{ERC2025}, about 2.3 billion a year, against 86 billion euros of higher-education research spending in the European Union in 2024 \citep{Eurostat2025}, gives $b$ near 0.03. A private foundation awarding on the order of \$100 million a year across several fields has $b$ at or below 0.01 in any one of them.


\clearpage
%----------------------------------------------------------------
% Supporting Information Appendix. At submission this block becomes the separate SI PDF;
% here it is compiled with the main text so that every \ref resolves. Order follows
% PNAS practice: SI Materials and Methods first, then numbered SI Text sections;
% SI figures and tables sit beside the text that uses them.
\section*{Supporting Information Appendix}
\setcounter{section}{0}\renewcommand{\thesection}{SI Text \arabic{section}}
\setcounter{figure}{0}\renewcommand{\thefigure}{S\arabic{figure}}
\setcounter{table}{0}\renewcommand{\thetable}{S\arabic{table}}
\setcounter{prop}{1} % the gap rule is restated as Proposition 1 in SI Text 1; the lottery result is Proposition 2

The SI Appendix holds what the main text cites without proof or states in one sentence: the model details, strategies, parameter ranges, the particle audit, and the settings behind every figure (SI Materials and Methods); the derivation of the gap rule and its comparison with the sequential optimum (\ref{app:gap-rule}); the lottery proposition (\ref{app:lottery}); four extended results with Figs.~\ref{fig:targeting-value} to~\ref{fig:concentration} (\ref{si:extended}); the robustness of the results to the production function (\ref{app:technology}); the sensitivity of the value of targeting and of review to the parameters the main sweeps hold fixed and to the signal model (\ref{si:sensitivity}); and the timing of funding (\ref{si:timing}). The table below maps each main-text statement to the section that supports it.

\begin{center}
\small
\begin{tabular}{@{}p{8.4cm}p{4.6cm}@{}}
\toprule
Main-text statement & Where in the SI \\
\midrule
Model, signals, strategies (Table~\ref{tab:strategies}), parameter ranges, copula, particles (Table~\ref{tab:particles}), scoring, settings behind every figure (Table~\ref{tab:settings}), code & SI Materials and Methods \\
The output-maximizing allocation of a round's budget funds the gap: $g_i^* = \max(cK_i - R_{i0}, 0)$ & \ref{app:gap-rule}, Proposition~1 \\
The funded set and every grant increase with the budget; the funding frontier & \ref{app:gap-rule}, Corollary~\ref{cor:frontier} \\
The value of targeting decreases in the budget and vanishes as the budget grows & \ref{app:gap-rule}, Corollary~\ref{cor:targeting-vanishes}; \ref{si:extended}, Fig.~\ref{fig:targeting-value} \\
Each heuristic can perform worse than uniform division & \ref{app:gap-rule}, Examples 1 and 2, Fig.~\ref{fig:heuristics} \\
The round-by-round gap rule is within about one percent of the sequential optimum & \ref{app:gap-rule}, Table~\ref{tab:dynopt} \\
Equal division outperforms any lottery over the same pool & \ref{app:lottery}, Proposition~\ref{prop:lottery} \\
Targeting's value relative to uniform funding, across budgets and fields & \ref{si:extended} \\
Overtrusting review loses more output than undertrusting it & \ref{si:extended}, Fig.~\ref{fig:overtrust} \\
Large grants reveal capability; the source of the shift of spending toward later rounds & \ref{si:extended}, Fig.~\ref{fig:depth-schedule} \\
The production function does not settle concentration & \ref{si:extended}, Fig.~\ref{fig:concentration} \\
The ordering of review's value by field holds from Cobb-Douglas to Leontief & \ref{app:technology}, Fig.~\ref{fig:signal-robustness} \\
The ordering of fields by the value of targeting and of review survives changes in mean capability, the capability-resource association, the resource tail, pool size, the resource signal, and the review-signal model; the inequality of the gaps organizes the results & \ref{si:sensitivity}, Tables~\ref{tab:sens-mean} to~\ref{tab:sens-contam}, Fig.~\ref{fig:gap-gini} \\
Planning the timing of spending adds little beside choosing recipients, and exists only under complementarity & \ref{si:timing}, Fig.~\ref{fig:p3}; \ref{app:technology}, Table~\ref{tab:timing-technology} \\
\bottomrule
\end{tabular}
\end{center}

% SI Materials and Methods: what the main text's Materials and Methods leaves out
% (interpretation, budget rule, signals, strategies, parameter ranges, scoring, the
% settings behind every figure, code). Built 2026-09-29 by merging the former SI Text 2
% (long-version Section 3) and SI Text 7 (Appendix D) and removing what Methods states.
\makeatletter
\section*{SI Materials and Methods}
\phantomsection\def\@currentlabel{SI Materials and Methods}\label{app:simulation}\label{sec:model}
\makeatother

\subsection*{Model details}
Output in each round is drawn from a Poisson distribution with rate $\lambda_i = AK_iR_i/(K_i + R_i)$, independently across researchers and rounds given the state; the units of output may be read as publications, findings, or units of scientific value. The production function is proportional to the harmonic mean of capability and resources. The harmonic mean makes the two inputs complements without making either an absolute requirement: a capable researcher with few resources still produces, by working more slowly or borrowing equipment, but no surplus of one input fully substitutes for a shortage of the other. \ref{app:technology} repeats the central comparisons for a family of production functions from Cobb-Douglas to Leontief.

Between rounds, capability grows by $\epsilon\,\Lambda(K_i, R_i)$ with $\Lambda = K_iR_i/(K_i + R_i)$: doing research builds expertise, the same bottleneck that limits output limits growth, and capability never declines. Growth follows expected output, so a researcher's accumulation depends on their capability and resources and not on luck in a given round. Resources do not accumulate. The baseline $R_{i0}$ recurs every round and represents the time and support a researcher has independent of the funder; a grant is consumed in the round it is given, and whatever lasts from it (skills, publications, standing) enters the model as capability.

The funder's total budget over the horizon is $b\,n\,\mathbb{E}[R_0]$, where $\mathbb{E}[R_0]$ is the expected baseline resources of one researcher, so at $b = 1$ the budget equals the community's expected baseline resources for one round. The budget does not grow with the horizon: a longer horizon varies the timing of spending, never its amount. Strategies that do not plan ahead spend $1/T$ of the budget each round; forward-looking strategies re-plan the entire remaining budget over the remaining horizon each round and execute the current round's allocation. For the timing runs (Fig.~3\textit{A} and Fig.~\ref{fig:depth-schedule}) the budget is scaled against capability instead, $b\,n\,\mathbb{E}[K_0]$, so that a community with few resources does not zero the budget. The simulations set the productivity constant $A = 1$; the code parameterizes the budget as $2\,b'\,n\,\mathbb{E}[R_0]$, so its $b'$ is half of the $b$ reported here.

The funder knows the distributions of capability and resources but observes no researcher's capability or resources directly. It observes each researcher's output as it accumulates. A resource signal, observed once at the start, gives a noisy estimate of each researcher's baseline resources (noise $\tau_R$), as one might obtain from institutional information. A review signal, drawn afresh in each round that a strategy uses it, gives a noisy estimate of capability (noise $\tau_K$), as one might obtain from proposal evaluation and expert judgment; the smaller $\tau_K$, the more informative the review. Formally, the resource signal is $R_{i0} + \eta_i$ with $\eta_i \sim N(0, \tau_R^2)$, and the review signal in round $t$ is $K_i + \xi_{it}$ with $\xi_{it} \sim N(0, \tau_K^2)$, independent across researchers and rounds; realized output is Poisson with mean $\lambda_i$. Bayesian strategies hold a posterior over each researcher's $(K_i, R_{i0})$, computed from these likelihoods by importance sampling with 200 particles drawn from the prior, update it after each round, and allocate each dollar to the researcher whose expected output, given the posterior, rises most with it, by exact water-filling on the posterior expected marginal values. The forward-looking funder is a certainty-equivalent receding-horizon planner: each round it chooses the current allocation and a plan for the remaining budget to maximize expected output over the remaining horizon, treating future output as equal to its expectation, and re-plans after observing the round's output; it does not value the information its own grants will generate.

\subsection*{Strategies}
Table~\ref{tab:strategies} lists the strategies compared. The Bayesian strategies differ in information, with or without the review signal, and in planning, myopic (optimizing the current round's expected output) or forward-looking (optimizing expected output over the remaining horizon). Three baselines complete the set: no funding, uniform funding, which splits each round's budget equally, and track-record funding, which allocates in proportion to the previous round's output. Seed grants and lotteries modify these strategies as described in the main text.

\begin{table}[t]
\centering
\footnotesize
\setlength{\tabcolsep}{4pt}
\begin{tabular*}{\textwidth}{@{\extracolsep{\fill}}lll@{}}
\toprule
Strategy & Information & Planning \\
\midrule
No funding & none & none \\
Uniform & none & none \\
Track record & track record & none \\
Myopic, without review & track record + resource signal & current round \\
Forward, without review & track record + resource signal & remaining horizon \\
Myopic, with review & track record + resource signal + review & current round \\
Forward, with review & track record + resource signal + review & remaining horizon \\
\bottomrule
\end{tabular*}
\caption{The funding strategies compared.}
\label{tab:strategies}
\end{table}

\subsection*{Parameter ranges and defaults}
The tail parameters $\alpha_K$ and $\alpha_R$ range from 1.3 to 3.5. Empirical distributions of research productivity have Pareto tail parameters near 1 to 2 \citep{Lotka1926, Nicholls1989}; % [check Nicholls range]
since productivity is an imperfect guide to capability (main text), the range extends well beyond those estimates toward more equal fields. The budget scale $b$ ranges from 0.2 to 2, from a funder whose budget is small beside the community's own resources to one able to relieve most researchers' bottlenecks, with larger budgets in the timing analyses. The compounding rate $\epsilon$ ranges from near zero, where funding raises only current output, to 0.85, where this round's funding substantially raises later capability. The review noise $\tau_K$ ranges from 0.05, nearly perfect review, to 20, nearly uninformative review, a range wide enough to contain the noise implied by empirical estimates of review's predictive accuracy. The correlation $\rho$ ranges from $-0.5$ to 0.8, and the horizon $T$ from one round to ten. Unless a figure's settings below say otherwise: $T = 2$; $n = 50$; capability and baseline resources drawn from Pareto distributions with scale 1 and $\alpha_K = \alpha_R = 2$; $\rho = 0$; $\epsilon = 0.1$; $b = 1$; $\tau_K = \tau_R = 1$; seed-grant fraction $0.25$ where seed grants are used.

\subsection*{Scoring and comparison}
Each strategy is scored by the sum of expected outputs along its run, which equals realized output averaged over repetitions of the same allocations; this removes simulation noise from comparisons without favoring any strategy. Every configuration is run on independently drawn populations, with the same random draws (populations, signals, and realized output) used for every strategy, so that comparisons between strategies are paired. Results are reported as shares of a named comparison, stated in each figure: most often the research output that funding adds (a strategy's expected output minus the expected output under no funding), and otherwise the records-only shortfall or the optimal allocation's gain. Standard errors, where shown, are across simulated populations.

\subsection*{Settings by figure}
Table~\ref{tab:settings} lists, for each figure, the settings that differ from the defaults.
\begin{table}[tp]
\centering
\small
\begin{tabular}{@{}p{2.3cm}p{10.0cm}@{}}
\toprule
Figure & Settings that differ from the defaults \\
\midrule
Fig.~2\textit{C} & Review noise $\tau_K$ from 0.05 to 20; $\alpha_K \in \{1.3, 2, 3.5\}$ with the Pareto scale set to $2(\alpha_K - 1)/\alpha_K$ so that mean capability is 2 in every field; 100 populations per point. The records-only shortfall relative to complete information is 45, 27, and 11 percent of the research output that complete-information funding adds, at the three tail parameters. \\
Fig.~2\textit{B} & $T = 20$; 50 populations. Records-only is the myopic funder without review; with review is the myopic funder with review. The shortfall is the funder's expected output below the complete-information optimum in that round, as a share of the optimum's gain over no funding. Under heavy-tailed capability ($\alpha_K = 1.3$, mean 2) the records-only shortfall falls from 61 to 16 percent against 7 to 3 percent with review, and the grant correlations run 0.07 to 0.43 against 0.92 in round one. \\
Fig.~\ref{fig:overtrust} & True noise $\tau_K = 3$; believed noise from 0.1 to 10; $\alpha_K \in \{1.3, 2\}$; 200 populations. Each negative value is within two standard errors of zero. \\
Fig.~\ref{fig:depth-schedule} & $T = 6$; community with few resources (Pareto scale for resources 0.001); $\alpha_K = 1.3$; no review signal ($\tau_K = 100$); $\epsilon = 10^{-4}$; budget scaled against capability; $b = 1$ (50 populations) and $b = 6$ (24 populations, standard error of each share at most 0.006). \\
Fig.~3\textit{A} & $T = 5$; forward-looking funder with review; budget scaled against capability; resource scales 0.001, 0.3, and 3; $\epsilon$ from $10^{-4}$ to 0.85; 200 populations. Five resource scales were run; the figure shows three, and the other two lie between them. \\
Fig.~3\textit{B} & $T$ from 2 to 10 (to 5 for $\epsilon \in \{0.05, 0.15, 0.55\}$); $\epsilon \in \{0.05, 0.15, 0.3, 0.55, 0.85\}$; 64 to 200 populations per point. The gain is the forward-looking funder's output over the myopic funder's, both with review; negative gains, all within noise of zero, are shown as zero. \\
Fig.~4\textit{A} & Review-informed forward-looking funder; seed-grant fraction from 0 to 0.75; heavy-tailed field ($\alpha_K = 1.3$) with reliable review ($\tau_K = 0.3$) at $b \in \{0.2, 1, 2\}$, and the default field at $b = 1$; 200 populations. The output lost is reported as a share of the review-informed funder's gain over no funding in the same setting. \\
Fig.~4\textit{B} & Single round; $n = 50$; a fifth of the population funded; review score $K_i$ plus Gaussian noise of standard deviation $\tau$; expected output computed exactly given each allocation; 2{,}000 populations per point (Monte Carlo standard error at most 0.002). Left: $\alpha_K = 1.3$, $b = 0.2$. Right: $\alpha_K = 3.5$, $b = 1$. \\
Fig.~\ref{fig:signal-robustness}, Table~\ref{tab:timing-technology} & Production functions with $\gamma \in \{0, -3\}$ and Leontief; 100 populations per point (50 for Leontief); $T = 2$ for the signal value and $T = 5$ for the timing table. These runs use the greedy allocator, the only one implemented for the general family (the harmonic-mean runs use exact water-filling), so their magnitudes are compared with one another and not with the main-text figures. Leontief allocations use a finer allocation step, with ties at the kink broken by fill order. \\
Fig.~\ref{fig:concentration} & Single round; review-informed funder; $\gamma$ from 0 to $-12$ and Leontief; three field-and-budget combinations ($\alpha_K = 1.3$, $b = 0.2$; $\alpha_K = 2$, $b = 2$; $\alpha_K = 3.5$, $b = 2$); 200 populations per point, 400 at the refinement points of the evenly spread curve. \\
Fig.~\ref{fig:gap-gini}, Tables~\ref{tab:sens-mean} to~\ref{tab:sens-contam} & $T = 2$; $b = 1$; myopic funders with and without review and the complete-information benchmark; $\rho \in \{-0.5, 0, 0.5, 0.8\}$; $\alpha_R \in \{1.3, 2, 3.5\}$; $n \in \{25, 50, 100, 200\}$; Pareto scale 1 or set to give mean capability 2; review score $K_i + \beta R_{i0}$ plus noise with $\beta \in \{0, 0.25, 0.5, 1\}$; 50 populations per setting, 100 in the noise sweeps. \\
\bottomrule
\end{tabular}
\caption{Simulation settings behind each figure.}
\label{tab:settings}
\end{table}

\subsection*{Code}
The model, the sweeps, and the scripts that produce every figure and number in the paper are in the project repository; the archived version will be cited at publication. % [check: repository URL and archived commit at publication]

% SI Text 1: derivation of the gap rule for constant-returns production functions (Proposition 1, Corollaries 1 and 2, Examples 1 and 2, Fig. S1).
\section{Derivation of the gap rule}\label{app:gap-rule}

\subsection*{The within-round allocation problem}

A population of $n$ researchers. Researcher $i$ has capability $K_i > 0$ and baseline resources $R_{i0} \geq 0$, both fixed within the round. A grant $g_i \geq 0$ adds to researcher $i$'s resources for the round, so that their expected output is $F(K_i,\, R_{i0} + g_i)$. We prove the gap rule for every production function of the form
\[
F(K, R) = A\,K\,h(R/K),
\]
where $A > 0$ is the productivity constant and $h\colon [0, \infty) \to [0, \infty)$ is continuous, continuously differentiable on $(0, \infty)$, strictly increasing, and strictly concave, so that $h'$ is positive and strictly decreasing. This is the class of production functions with constant returns to scale in $(K, R)$ and diminishing returns to resources. The harmonic mean of the main text is the instance $h(x) = x/(1 + x)$, for which $F(K, R) = AKR/(K + R)$; every member of the power-mean family of \ref{app:technology} with exponent $\gamma < 1$ is an instance, with $h(x) = \big((1 + x^{\gamma})/2\big)^{1/\gamma}$. Write
\[
\varphi_i(g) = A\,K_i\,h\!\left(\frac{R_{i0} + g}{K_i}\right)
\]
for researcher $i$'s expected output as a function of their grant alone; for the harmonic mean, $\varphi_i(g) = A K_i (R_{i0} + g)/(K_i + R_{i0} + g)$.

Throughout, the round's budget is $B > 0$. Nothing else about the model is used, so the result holds for every realized population and every round's budget; across rounds, the constant $c$ and the funded set are recomputed from the new population state.

The funder observes $(K_i, R_{i0})_{i=1}^n$ and divides the budget so as to maximize expected output:
\[
(\mathrm{P}) \qquad \max_{g_1, \ldots, g_n \geq 0} \; \sum_{i=1}^n \varphi_i(g_i) \quad \text{subject to} \quad \sum_{i=1}^n g_i \leq B.
\]

\subsection*{Three lemmas}

\begin{lemma}[the return to a grant]\label{lem:return}
For each $i$, on $g \in [0, \infty)$: (i) $\varphi_i$ is strictly increasing, with $\varphi_i'(g) = A\,h'\!\big((R_{i0} + g)/K_i\big) > 0$; (ii) $\varphi_i$ is strictly concave, since $\varphi_i'$ is strictly decreasing in $g$; (iii) $\varphi_i(g) \uparrow A K_i \bar h$ as $g \to \infty$, where $\bar h = \lim_{x \to \infty} h(x) \leq \infty$. For the harmonic mean, $\varphi_i'(g) = A\,K_i^2/(K_i + R_{i0} + g)^2$, $\varphi_i''(g) = -\,2A\,K_i^2/(K_i + R_{i0} + g)^3$, $\varphi_i'(g) < A$ whenever $R_{i0} + g > 0$, and $\bar h = 1$, so expected output never exceeds $AK_i$.
\end{lemma}

\begin{proof}
Differentiating $\varphi_i(g) = A K_i h((R_{i0} + g)/K_i)$ in $g$ gives (i), and (ii) follows because $h'$ is strictly decreasing and its argument increases with $g$. For (iii), $h$ is increasing, so $h(x) \uparrow \bar h$ as $x \to \infty$. For the harmonic mean, write $S = K_i + R_{i0} + g$; then $\varphi_i(g) = AK_i - AK_i^2/S$, and the stated derivatives follow (note $dS/dg = 1$); $\varphi_i'(g) = AK_i^2/S^2 < A$ exactly when $S > K_i$, that is, when $R_{i0} + g > 0$.
\end{proof}

\Cref{lem:return} says that the marginal value of a dollar to researcher $i$ depends on capability and resources only through their ratio, decreases as that ratio rises, and so increases with capability and decreases with the resources they already hold; that each dollar granted lowers the value of the next; and that, when $h$ is bounded, expected output never exceeds a ceiling set by capability.

\begin{lemma}[existence, uniqueness, exhaustion]\label{lem:existence}
Problem $(\mathrm{P})$ has exactly one solution $g^* = (g_1^*, \ldots, g_n^*)$, and it exhausts the budget: $\sum_i g_i^* = B$.
\end{lemma}
\begin{proof}
The feasible set is compact and convex and the objective $f(g) = \sum_i \varphi_i(g_i)$ is continuous and, by \cref{lem:return}(ii), strictly concave, so a maximizer exists and is unique. If $\sum_i g_i^* < B$, increasing any one $g_i^*$ slightly stays feasible and raises $f$ (\cref{lem:return}(i)), contradicting optimality.
\end{proof}
\begin{lemma}[equal marginal value]\label{lem:emv}
A feasible allocation $g$ with $\sum_i g_i = B$ solves $(\mathrm{P})$ if and only if there is a $\nu > 0$ such that
\[
\varphi_i'(g_i) = \nu \;\text{ for every } i \text{ with } g_i > 0, \qquad \varphi_i'(0) \leq \nu \;\text{ for every } i \text{ with } g_i = 0,
\]
where $\varphi_i'(0)$ may be $+\infty$ (as it is when $R_{i0} = 0$ and $h'(0^+) = \infty$), in which case researcher $i$ is funded.
\end{lemma}

\begin{proof}
\emph{Necessity.} Let $g^*$ solve $(\mathrm{P})$. Since $B > 0$ is exhausted (\cref{lem:existence}), some researcher is funded: $g_i^* > 0$ for some $i$. Fix any such $i$ and any $j \neq i$, and for $0 < \delta \leq g_i^*$ consider the transfer that replaces $g_i^*$ by $g_i^* - \delta$ and $g_j^*$ by $g_j^* + \delta$, leaving the rest unchanged. The result is feasible, so optimality gives
\[
\varphi_j(g_j^* + \delta) - \varphi_j(g_j^*) \;\leq\; \varphi_i(g_i^*) - \varphi_i(g_i^* - \delta).
\]
Dividing by $\delta$ and letting $\delta \downarrow 0$ yields $\varphi_j'(g_j^*) \leq \varphi_i'(g_i^*)$. If $j$ is also funded, the same argument with $i$ and $j$ exchanged gives the reverse inequality, so all funded researchers share a common marginal value; call it $\nu = \varphi_i'(g_i^*) > 0$ (positivity by \cref{lem:return}(i)). For unfunded $j$, the displayed inequality reads $\varphi_j'(0) \leq \nu$.

\emph{Sufficiency.} Let $g$ satisfy the condition with multiplier $\nu$, and let $\tilde g$ be any feasible allocation. Concavity of $\varphi_i$ gives, for every $i$, $\varphi_i(\tilde g_i) \leq \varphi_i(g_i) + \varphi_i'(g_i)\,(\tilde g_i - g_i)$. For funded $i$, $\varphi_i'(g_i) = \nu$, so $\varphi_i'(g_i)(\tilde g_i - g_i) = \nu\,(\tilde g_i - g_i)$. For unfunded $i$, $g_i = 0$ and $\tilde g_i - g_i = \tilde g_i \geq 0$, so $\varphi_i'(0)\,(\tilde g_i - g_i) \leq \nu\,(\tilde g_i - g_i)$. Summing over $i$:
\[
f(\tilde g) \;\leq\; f(g) + \nu \sum_i (\tilde g_i - g_i) \;=\; f(g) + \nu\Big(\sum_i \tilde g_i - B\Big) \;\leq\; f(g),
\]
since $\tilde g$ is feasible and $\nu > 0$. So $g$ solves $(\mathrm{P})$.
\end{proof}

\Cref{lem:emv} is the formal content of the sentence in the main text: an allocation is optimal exactly when no dollar can be moved to a researcher who can produce more value with it.

\subsection*{The proposition}

\begin{proprestate}[the gap rule]
Let $F(K, R) = A\,K\,h(R/K)$ with $h$ increasing and strictly concave. Then problem $(\mathrm{P})$ has the unique solution
\[
g_i^* = \max(c\,K_i - R_{i0},\, 0), \qquad i = 1, \ldots, n,
\]
where $c > 0$ is the unique constant satisfying the budget identity $\sum_{i=1}^n \max(c\,K_i - R_{i0},\, 0) = B$. Equivalently, $c = (h')^{-1}(\nu/A)$, where $\nu$ is the common marginal value of \cref{lem:emv}; for the harmonic mean, $c = \sqrt{A/\nu} - 1$ with $\nu \in (0, A)$.
\end{proprestate}

\begin{proof}
Let $g^*$ be the unique solution (\cref{lem:existence}) and $\nu$ the multiplier that \cref{lem:emv} associates with it. The proof proceeds in four steps.

\emph{Step 1: $c = (h')^{-1}(\nu/A)$ is well defined and positive.} Some researcher $i$ is funded, and for them \cref{lem:emv} gives $\nu = \varphi_i'(g_i^*) = A\,h'\big((R_{i0} + g_i^*)/K_i\big)$, so $\nu/A$ lies in the range of $h'$; since $h'$ is strictly decreasing, its inverse is defined there, and $c = (h')^{-1}(\nu/A) = (R_{i0} + g_i^*)/K_i > 0$ because $g_i^* > 0$. For the harmonic mean, $h'(x) = 1/(1 + x)^2$, so $(h')^{-1}(y) = 1/\sqrt{y} - 1$ and $c = \sqrt{A/\nu} - 1$; positivity of $c$ is then $\nu < A$, which \cref{lem:return} gives directly.

\emph{Step 2: funded researchers receive exactly their gap.} Let $g_i^* > 0$. By \cref{lem:emv}, $A\,h'\big((R_{i0} + g_i^*)/K_i\big) = \nu$, so $(R_{i0} + g_i^*)/K_i = (h')^{-1}(\nu/A) = c$ and $g_i^* = c\,K_i - R_{i0}$. Since $g_i^* > 0$, this equals $\max(c\,K_i - R_{i0},\, 0)$. Every funded researcher is thus brought to the same ratio of resources to capability, $c$.

\emph{Step 3: unfunded researchers have no gap.} Let $g_i^* = 0$. By \cref{lem:emv}, $A\,h'(R_{i0}/K_i) = \varphi_i'(0) \leq \nu = A\,h'(c)$, and since $h'$ is strictly decreasing, $R_{i0}/K_i \geq c$, that is, $R_{i0} \geq c\,K_i$. Then $\max(c\,K_i - R_{i0},\, 0) = 0 = g_i^*$.

Steps 2 and 3 establish $g_i^* = \max(c\,K_i - R_{i0}, 0)$ for every $i$, and the budget identity is exhaustion (\cref{lem:existence}). It remains to show $c$ is the only constant satisfying that identity.

\emph{Step 4: the budget identity pins $c$ uniquely.} Define $G(c) = \sum_{i=1}^n \max(c\,K_i - R_{i0},\, 0)$ for $c \geq 0$. $G$ is continuous and nondecreasing, with $G(0) = 0$ and $G(c) \to \infty$ as $c \to \infty$. Moreover $G$ is strictly increasing wherever it is positive: if $G(c) > 0$, then $c\,K_j > R_{j0}$ for some $j$, and for any $c' > c$, $G(c') \geq G(c) + K_j\,(c' - c) > G(c)$. Now suppose $c_1 < c_2$ both satisfy $G(c) = B$. Since $B > 0$, $G(c_1) > 0$, so $G(c_2) > G(c_1) = B$, a contradiction.
\end{proof}

The proof uses only that the marginal value of a dollar depends on $K$ and $R$ through the ratio $R/K$ and falls as that ratio rises. Constant returns to scale give the first property and diminishing returns to resources the second; the harmonic form contributes nothing beyond the closed form of $c$. The constant $c$ is the ratio of resources to capability at which the optimal allocation leaves every funded researcher, and so also the rate at which the rule converts a unit of capability into a target level of resources.

\subsection*{Consequences}

\begin{corollary}[funding frontier and budget comparative statics]\label{cor:frontier}
(i) Researcher $i$ is funded, $g_i^* > 0$, exactly when $K_i > R_{i0}/c$. (ii) As a function of the budget $B$, the constant $c(B)$ is continuous and strictly increasing; hence the funded set is nondecreasing in $B$, every funded researcher's grant $c(B)\,K_i - R_{i0}$ is strictly increasing in $B$, and the common marginal value $\nu(B) = A\,h'(c(B))$ is strictly decreasing in $B$. Both parts follow from the formula and Step 4.
\end{corollary}
\begin{corollary}[the value of targeting vanishes in the budget]\label{cor:targeting-vanishes}
Let $f^*(B)$ denote the optimal value of $(\mathrm{P})$ and $f^u(B) = \sum_i \varphi_i(B/n)$ the value of uniform funding. Then $0 \leq f^*(B) - f^u(B)$ for every $B$. If $h$ is bounded, $\bar h < \infty$, then $f^*(B) - f^u(B) \to 0$ as $B \to \infty$.
\end{corollary}

\begin{proof}
The first inequality holds because uniform funding is feasible. For the second: by \cref{lem:return}(iii), $\varphi_i(g) < A K_i \bar h$ for every $g$, so $f^*(B) \leq A\bar h\sum_i K_i$; and $f^u(B) = \sum_i \varphi_i(B/n) \to A\bar h\sum_i K_i$ as $B \to \infty$. Hence $0 \leq f^*(B) - f^u(B) \leq A\bar h\sum_i K_i - f^u(B) \to 0$.
\end{proof}

Two remarks. First, the corollary says nothing about monotonicity: $f^*(B) - f^u(B)$ is zero at $B = 0$ and positive for small $B > 0$ whenever researchers differ, so it cannot decrease throughout. The decrease reported in the main text is a numerical finding about the ratio $(f^* - f^u)/(f^u - f^0)$, with $f^0$ the output under no funding: it fell monotonically in the budget, on budget scales from $0.02$ to $3$, in every one of 300 populations of 50 researchers drawn from the default distributions, and in the population of Fig.~\ref{fig:targeting-value}. Since uniform funding's own gain over no funding does not vanish (it approaches $A\bar h\sum_i K_i - \sum_i \varphi_i(0) > 0$), the ratio vanishes with the difference. Second, boundedness of $h$ is what makes targeting's value vanish: in the power-mean family, $h$ is bounded exactly when $\gamma < 0$, with $\bar h = 2^{-1/\gamma}$ (so $\bar h = 1$ for the harmonic mean), whereas under Cobb-Douglas production ($\gamma \to 0$, $h(x) = \sqrt{x}$) output has no ceiling set by capability and $f^*(B) - f^u(B)$ grows without bound in $B$ for any population with unequal capabilities. The Leontief form $h(x) = \min(1, x)$ is the limit $\gamma \to -\infty$ of the family; it is not strictly concave, so the proposition does not apply to it directly, and it is treated numerically in \ref{app:technology}.

\begin{example}[funding the track record can produce less than uniform funding]\label{ex:track-record}
Two researchers, harmonic production, $A = 1$, budget $B = 2$. Researcher 1: $K_1 = 8$, $R_{10} = 8$ (productive and well resourced; expected output $\varphi_1(0) = 4$). Researcher 2: $K_2 = 4$, $R_{20} = 1$ (capable but with few resources; $\varphi_2(0) = 4/5$). Allocating in proportion to expected output gives $g = (5/3,\, 1/3)$ and total output $285/53 \approx 5.38$. Uniform funding, $g = (1, 1)$, gives $284/51 \approx 5.57$. The gap rule ($c = 3/4$) gives $g^* = (0, 2)$ and $40/7 \approx 5.71$.
\end{example}

\begin{example}[funding the under-resourced can produce less than uniform funding]\label{ex:scarcity}
Two researchers, harmonic production, $A = 1$, budget $B = 2$. Researcher 1: $K_1 = 10$, $R_{10} = 5$. Researcher 2: $K_2 = 1/2$, $R_{20} = 0$ (the scarcest resources in the population). Directing the budget to the researcher with the fewest resources gives total output $56/15 \approx 3.73$. Uniform funding gives $49/12 \approx 4.08$. The gap rule ($c = 2/3$) gives $g^* = (5/3,\, 1/3)$ and $21/5 = 4.2$.
\end{example}

% Figure (Appendix A): the two intuitive schemes on the population of Figure 1.
% Same 75 researchers as Figure 1 (fig1h-all.dat). Budget B_1 = 8.52, the total of the
% grants at the smaller budget in Figure 1 (c_1 = 2/3), divided equally among nine
% researchers (x = 0.947 each). Left: the nine with the highest expected output
% K R/(K + R) (fig1h-track.dat). Right: the nine with the fewest baseline resources
% (fig1h-under.dat). Expected output over no funding on this population (A = 1):
% gap rule +4.01, track record +1.98, under-resourced +2.09, uniform +1.99.
\begin{figure}[tp]
\centering
\begin{tikzpicture}
\begin{axis}[
  name=left,
  width=0.46\textwidth, height=5.2cm,
  axis lines*=left,
  xmin=0, xmax=8.6, ymin=0, ymax=7.6,
  xtick=\empty, ytick=\empty,
  x axis line style={-{Stealth[length=4pt]}, line width=0.4pt},
  y axis line style={-{Stealth[length=4pt]}, line width=0.4pt},
  xlabel={resources $R$},
  ylabel={capability $K$},
  ylabel style={font=\small}, xlabel style={font=\small},
  title={\small funding the track record},
  title style={yshift=-4pt},
  axis line style={line width=0.4pt},
  clip=false,
]
\addplot[only marks, mark=o, mark size=1.3pt, gray, line width=0.4pt] table[x=R, y=K] {fig1h-all.dat};
\addplot[quiver={u=\thisrow{u}, v=\thisrow{v}}, -{Stealth[length=3pt]}, gray, line width=0.5pt] table {fig1h-track.dat};
\addplot[only marks, mark=*, mark size=1.5pt, black] table[x=R, y=K] {fig1h-track.dat};
\end{axis}
\begin{axis}[
  at={(left.outer east)}, anchor=outer west, xshift=4pt,
  width=0.46\textwidth, height=5.2cm,
  axis lines*=left,
  xmin=0, xmax=8.6, ymin=0, ymax=7.6,
  xtick=\empty, ytick=\empty,
  x axis line style={-{Stealth[length=4pt]}, line width=0.4pt},
  y axis line style={-{Stealth[length=4pt]}, line width=0.4pt},
  xlabel={resources $R$},
  xlabel style={font=\small},
  title={\small funding the under-resourced},
  title style={yshift=-4pt},
  axis line style={line width=0.4pt},
  clip=false,
]
\addplot[only marks, mark=o, mark size=1.3pt, gray, line width=0.4pt] table[x=R, y=K] {fig1h-all.dat};
\addplot[quiver={u=\thisrow{u}, v=\thisrow{v}}, -{Stealth[length=3pt]}, gray, line width=0.5pt] table {fig1h-under.dat};
\addplot[only marks, mark=*, mark size=1.5pt, black] table[x=R, y=K] {fig1h-under.dat};
\end{axis}
\end{tikzpicture}
\caption{Two heuristics on the population of Fig.~\ref{fig:p1}. The budget of that figure's smaller frontier is divided equally among nine researchers. Left: funding the track record gives the grants (arrows) to the nine researchers with the highest expected research output, most of them well resourced and of middling capability. Right: funding the under-resourced gives the grants to the nine researchers with the fewest resources, most of them of low capability. On this population the gap rule adds 4.0 units of expected research output over no funding; funding the track record adds 2.0, funding the under-resourced 2.1, and dividing the budget uniformly across all 75 researchers 2.0.}
\label{fig:heuristics}
\end{figure}


% SI Text 2: lotteries and equal division (Proposition 2).
\section{Lotteries and equal division}\label{app:lottery}

Fix a round, a budget $B > 0$, and a pool $P \subseteq \{1, \ldots, n\}$ of $m$ researchers. The \emph{equal division} over $P$ grants each researcher in $P$ the amount $B/m$. A \emph{lottery} over $P$ with $k \leq m$ winners selects $k$ researchers from $P$ uniformly at random and grants each winner $B/k$. Expected output is evaluated as in \ref{app:gap-rule}: given an allocation $g$, it is $\sum_i \varphi_i(g_i)$, and for a random allocation, the expectation is taken over the draw as well.

\begin{prop}[lotteries are outperformed by equal division]\label{prop:lottery}
For every pool $P$, budget $B$, and winner count $k \leq m$, the lottery's expected output does not exceed the equal division's, with equality only if $k = m$. In particular, a lottery over the whole population, at any winner count, produces at most uniform funding's expected output.
\end{prop}

\begin{proof}
Researchers outside $P$ receive nothing under either scheme. For $i \in P$, the lottery's grant to $i$ is a random variable $\tilde g_i$ equal to $B/k$ with probability $k/m$ and $0$ otherwise, so $\mathbb{E}[\tilde g_i] = B/m$. By linearity, the lottery's expected output is $\sum_{i \notin P} \varphi_i(0) + \sum_{i \in P} \mathbb{E}[\varphi_i(\tilde g_i)]$. Since $\varphi_i$ is strictly concave (Lemma~\ref{lem:return}(ii)), Jensen's inequality gives $\mathbb{E}[\varphi_i(\tilde g_i)] \leq \varphi_i(\mathbb{E}[\tilde g_i]) = \varphi_i(B/m)$ for every $i \in P$, with equality only if $\tilde g_i$ is degenerate, that is, $k = m$. Summing over $i$ yields the claim; taking $P = \{1, \ldots, n\}$ yields the comparison with uniform funding.
\end{proof}

\begin{remark}[what the proposition does and does not cover]
The proposition compares expected output within a fixed pool. It measures the effect of the randomization alone: whatever pool review selects, replacing the lottery within that pool by the equal division of the same money weakly increases expected output. The proposition says nothing about the choice of pool, about the review ranking within the pool that the equal division also ignores, or about anything a lottery saves outside the model (review costs, proposal effort, bias).
\end{remark}

\begin{remark}[minimum viable grant size]
The proposition needs divisible grants and diminishing returns at every grant size. Suppose instead that a project produces nothing unless its grant reaches a minimum $m$, and that $B < m\,|P|$ for a pool $P$. Equal division then gives every researcher in $P$ less than $m$ and produces no output, while a lottery with $k = \lfloor B/m \rfloor$ winners, each receiving $B/k \geq m$, produces the output of $k$ funded projects. Where minimum project scale binds, the comparison of the proposition reverses, and this qualification applies wherever the main text compares lotteries with equal division.
\end{remark}


% SI Text 5: extended results. Full treatment of results the main text states in one
% sentence; drawn from the long version's Sections 4 to 7 (paragraphs and footnotes),
% with Figs. S1 to S4.
\section{Extended results}\label{si:extended}

This section gives the full treatment of results that the main text states in one sentence. Each subsection opens with the main-text statement it supports.

\subsection*{The value of targeting decreases in the budget}
\emph{Main text:} ``The value of targeting, the expected output the optimal allocation adds over equal division, decreases in the budget and vanishes as the budget grows.''

How much targeting matters depends on the budget. A researcher's expected output increases with every dollar of resources, but with diminishing marginal returns, and it never exceeds $AK_i$: expected output is bounded above by a quantity set by capability, which it approaches but does not attain. When the budget is large, optimal and uniform funding therefore differ little in the output they produce: under either, every researcher's output is near its upper bound, and the difference between the two vanishes as the budget increases (Corollary~\ref{cor:targeting-vanishes}, \ref{app:gap-rule}). When the budget is small, dollars are scarce, researchers differ in the marginal value of a dollar, and output depends on where each dollar goes. Targeting matters where the budget is small relative to researchers' resources and little where the budget is large (Fig.~\ref{fig:targeting-value}). This dependence on the budget recurs later: it sets how much output is lost when seed grants and lotteries override the optimal allocation, and it determines when optimal funding concentrates rather than spreads (both in the main text).

% Figure 2: the value of targeting decreases in the budget.
% Computed exactly from Proposition 1 (no simulation): population of n = 75 with
% K, R0 ~ Pareto(alpha = 2, min 1), rho = 0, A = 1, seed 659 (numpy default_rng);
% value of targeting = (f_opt - f_unif) / (f_unif - f_0), in percent;
% b = budget / total baseline resources. Data: fig2-curve.dat, generated by
% verify_gap_rule.py's companion snippet (see state.md, 2026-09-05).
\begin{figure}[tp]
\centering
\begin{tikzpicture}
\begin{axis}[
  width=0.82\textwidth, height=5.4cm,
  axis lines*=left,
  xmin=0, xmax=2, ymin=0, ymax=175,
  xtick={0, 0.5, 1, 1.5, 2},
  ytick={0, 50, 100, 150},
  xlabel={budget scale $b$},
  ylabel={value of targeting (\%)},
  ylabel style={font=\small}, xlabel style={font=\small},
  tick label style={font=\small},
  axis line style={line width=0.4pt},
  tick style={line width=0.4pt, black},
  clip=false,
]
\addplot[black, line width=0.9pt, smooth] table {fig2-curve.dat};
\end{axis}
\end{tikzpicture}
\caption{The value of targeting against the budget. The optimal allocation's gain over uniform funding, as a fraction of uniform funding's gain over no funding, as the budget scale $b$ grows, for one population drawn from the default distributions. Corollary~\ref{cor:targeting-vanishes} proves that the value vanishes as the budget grows; that it falls throughout is a numerical finding (\ref{app:gap-rule}).}
\label{fig:targeting-value}
\end{figure}


\subsection*{The value of targeting relative to uniform funding}
\emph{Main text:} ``For a funder allocating with records and review, targeting adds more research output than uniform funding's entire gain over no funding in a heavy-tailed field with reliable review and a tight budget, and about a quarter of that gain in the default field at the largest budget we sweep.''

The review-informed funder's gain over uniform funding, as a fraction of uniform funding's gain over no funding, decreases from 1.17 to 0.72 across budget scales $b = 0.2$ to $2$ with heavy-tailed capability ($\alpha_K = 1.3$) and reliable review ($\tau_K = 0.3$), and from 0.74 to 0.25 in the default field with the default signal ($\alpha_K = 2$, $\tau_K = 1$). Two rounds; 200 simulated populations per point (the populations of Fig.~4\textit{A}). Giving up targeting loses little output under the same conditions that make review worth little: where capability is evenly spread, the budget is large, or review is uninformative.

\subsection*{Overtrusting review}
\emph{Main text:} ``A funder that overtrusts review, treating a noisy signal as reliable, can do worse than one that ignores review altogether, and overtrusting review loses more output than undertrusting it.''

One condition qualifies these results: the funder must weight review at its true informativeness. A funder that overtrusts review, treating a noisy signal as reliable, can do worse than one that ignores review altogether: with moderately spread capability, overtrust turns review's value negative; with heavy-tailed capability it reduces the value without reversing it (Fig.~\ref{fig:overtrust}).\footnote{A funder that overtrusts tenfold (true $\tau_K = 3$, belief $0.3$) loses more than review's full value at the default capability distribution; a funder that undertrusts tenfold (true $\tau_K = 0.3$, belief $3$) retains over half of it.} Overtrusting review loses more output than undertrusting it.

% Figure 6 (S6): the value of review under misweighted trust.
% y = review's value (S5 - S4) as percent of no-funding output; x = the funder's
% believed noise (log scale); true noise tau_K = 3 (dotted rule). Two capability
% distributions. Error bars: one standard error, 200 seeds per cell.
% Data: fig6-overtrust.dat, from sweep_results/D_misspecified_trust (hash 21b0d9a).
\begin{figure}[tp]
\centering
\begin{tikzpicture}
\begin{axis}[
  width=0.82\textwidth, height=6.0cm,
  axis lines*=left,
  xmode=log,
  xmin=0.1, xmax=10, ymin=-5, ymax=45,
  xtick={0.1, 0.3, 1, 3, 10},
  xticklabels={0.1, 0.3, 1, 3, 10},
  ytick={0, 10, 20, 30, 40},
  xlabel={believed noise $\tau_K$},
  ylabel={value of review (\% of the gain from funding)},
  ylabel style={font=\small}, xlabel style={font=\small},
  tick label style={font=\small},
  axis line style={line width=0.4pt},
  tick style={line width=0.4pt, black},
  clip=false,
]
\addplot[gray, dashed, line width=0.4pt, domain=0.1:10] {0};
\draw[gray, dotted, line width=0.5pt] (axis cs:3,-2) -- (axis cs:3,20);
\node[font=\small, gray, anchor=west] at (axis cs:3.15,18) {true noise};
\addplot[black, line width=0.7pt, error bars/.cd, y dir=both, y explicit]
  table[x=belief, y=p13, y error=s13] {fig6-overtrust.dat};
\addplot[black, line width=0.7pt, error bars/.cd, y dir=both, y explicit]
  table[x=belief, y=p20, y error=s20] {fig6-overtrust.dat};
\node[font=\small, anchor=west] at (axis cs:11,25.5) {$\alpha_K = 1.3$};
\node[font=\small, anchor=west] at (axis cs:11,3.9)  {$\alpha_K = 2$};
\end{axis}
\end{tikzpicture}
\caption{Overtrusting review loses more research output than undertrusting it. The value of review as a percent of the research output that a correctly calibrated review-informed funder adds over no funding, as the funder's believed signal noise varies while the true noise is $\tau_K = 3$ (dotted rule); bars are one standard error. At the default capability distribution, overtrust takes the value below zero. Settings: \ref{app:simulation}.}
\label{fig:overtrust}
\end{figure}


\subsection*{Grant size, learning, and the schedule}
\emph{Main text:} ``The shift of spending toward later rounds comes from the output researchers produce whether or not they are funded, which compounds and reveals capability while the funder waits.''

How much a funder can learn from its own grants depends on their size. A researcher on a small grant produces output nearly proportional to the grant and nearly independent of capability, so a funder whose grants are small learns almost nothing it can exploit later: in a field whose researchers have few resources, with no review signal, such a funder gains almost nothing over uniform funding. Large grants change this. As grants approach the scale of capability itself, output separates the capable from the rest, and the funder recovers most of the value of discrimination, on a schedule that spends a small share early, so that output can be observed, and the rest once the funder has learned who is capable (Fig.~\ref{fig:depth-schedule}).\footnote{In a community whose researchers have few resources, with heavy-tailed capability and no review signal, over six rounds, the Bayesian funder's gain over uniform funding increases from a third of one percent of uniform funding's output at the default grant size to nine percent at six times that size and ten percent at twelve times, roughly the gain a reliable review signal delivers at the default size.}

The shift of spending toward later rounds comes from the output researchers produce whether or not they are funded: that output compounds and reveals capability while the funder waits. Growth from grant-added output alone would shift the schedule slightly toward earlier rounds, since earlier grants have longer to compound.\footnote{With growth from unfunded output switched off, the schedule's center of mass is 0.48 to 0.49, below even; with growth from grant-added output switched off, 0.53 to 0.68 (five rounds).}

% Figure (S7): the optimal schedule at standard and deep funding.
% y = share of the budget spent in each round by the forward funder (S8); T = 6;
% poverty corner (r_min = 0.001), heavy tail (alpha_K = 1.3), no review signal
% (tau_K = 100), negligible compounding (eps = 1e-4), budget scaled against
% capability (budget_ref = "K"). b = 0.5: 50 seeds; b = 3: 24 seeds (max SE 0.006);
% round-1 share at b = 3 matches the canonical 200-seed run (0.104 vs 0.110).
% Data: fig7-schedule.dat, generated by for-claude/verify_s7.R (schedule addendum).
\begin{figure}[tp]
\centering
\begin{tikzpicture}
\begin{axis}[
  width=0.66\textwidth, height=5.4cm,
  axis lines*=left,
  xmin=1, xmax=6, ymin=0.08, ymax=0.20,
  xtick={1, 2, 3, 4, 5, 6},
  ytick={0.10, 0.15, 0.167, 0.20},
  yticklabels={0.10, 0.15, even, 0.20},
  xlabel={round},
  ylabel={share of budget},
  ylabel style={font=\small}, xlabel style={font=\small},
  tick label style={font=\small},
  axis line style={line width=0.4pt},
  tick style={line width=0.4pt, black},
  clip=false,
]
\addplot[black, line width=0.7pt, mark=*, mark size=1.2pt] table[x=round, y=b05] {fig7-schedule.dat};
\addplot[black, line width=0.7pt, mark=o, mark size=1.4pt] table[x=round, y=b3] {fig7-schedule.dat};
\node[font=\small, anchor=west] at (axis cs:6.15,0.167) {default grant size ($b = 1$)};
\node[font=\small, anchor=west] at (axis cs:6.15,0.183) {large grants ($b = 6$)};
\end{axis}
\end{tikzpicture}
\caption{The forward-looking funder's spending by round at two grant sizes. The share of the budget spent in each of six rounds, in a field with heavy-tailed capability, few baseline resources, and no review, at the default grant size (filled) and at six times that size (open). Settings: \ref{app:simulation}.}
\label{fig:depth-schedule}
\end{figure}


\subsection*{Concentration and the production function}
\emph{Main text:} ``Stronger complementarity between capability and resources concentrates the review-informed funder's grants when the budget is tight in a heavy-tailed field, spreads them when the budget is ample, and has a non-monotone effect where capability is evenly spread.''

Seed grants and lotteries both spread funds more widely than targeting does, and both raise the older question of whether funding should be concentrated among a few researchers or spread across many. One might hope that the production function settles this question: the harder it is to substitute resources for capability, the more the money should go to the few researchers who can use it. To test this, we embed our production function in a family indexed by an exponent $\gamma$, from Cobb-Douglas ($\gamma = 0$), under which a shortfall in either input can be offset by more of the other, through our form ($\gamma = -1$), to Leontief, under which output depends only on the scarcer input; the harder it is to substitute one input for the other, the stronger the complementarity (\ref{app:technology}). The production function does not settle the question (Fig.~\ref{fig:concentration}). With a tight budget in a heavy-tailed field, stronger complementarity does concentrate the review-informed funder's grants. With an ample budget, stronger complementarity spreads them, a pattern consistent with a funder that, once the budget allows, brings every researcher's resources up toward the level their capability can use. Where capability is evenly spread and the budget ample, the relation is not even monotone: concentration peaks at an intermediate complementarity and declines toward both ends.\footnote{Gini coefficients of the review-informed funder's grants, from Cobb-Douglas to Leontief: 0.92 to 0.95 in a heavy-tailed field with a tight budget; 0.35 to 0.21 in the default field with an ample budget; 0.63 to 0.45 in a heavy-tailed field with an ample budget. In the evenly spread field with an ample budget the Gini coefficient is 0.16 at Cobb-Douglas, peaks at 0.196 near $\gamma = -6$, and falls to 0.185 at $\gamma = -12$ (3.5 standard errors above the $\gamma = -12$ level).} How widely funds should be spread is set jointly by the production function, the budget, and the field's capability distribution; no one of the three decides it alone.

% Figure (S8): how concentrated the informed allocation is, across the family.
% y = Gini coefficient of the review-informed funder's grants (single round);
% x = CES exponent gamma (0 = Cobb-Douglas, -1 = our harmonic form, limit =
% Leontief, plotted as detached marks at the left edge). Three field-budget pairs.
% Ground: sigma_tierB_concentration (+ leontief endpoint, + 400-seed refinement for
% the evenly spread curve), 200 seeds per cell, re-derived in-session (verify_s9.R).
% Leontief cells carry the kink caveat (ties broken by fill order, n_steps = 800).
\begin{figure}[tp]
\centering
\begin{tikzpicture}
\begin{axis}[
  width=0.82\textwidth, height=6.0cm,
  axis lines*=left,
  xmin=-14.8, xmax=0, ymin=-0.16, ymax=0.06,
  xtick={-14, -12, -6, -3, -1, 0},
  xticklabels={Leontief, $-12$, $-6$, $-3$, $-1$, $0$},
  ytick={-0.15, -0.10, -0.05, 0, 0.05},
  yticklabels={$-0.15$, $-0.10$, $-0.05$, 0, 0.05},
  xlabel={CES exponent $\gamma$ (complements $\leftarrow$, substitutable $\rightarrow$)},
  ylabel={change in grant concentration (Gini)},
  ylabel style={font=\small}, xlabel style={font=\small},
  tick label style={font=\small},
  axis line style={line width=0.4pt},
  tick style={line width=0.4pt, black},
  clip=false,
]
\draw[gray, dotted, line width=0.5pt] (axis cs:-1,-0.16) -- (axis cs:-1,0.06);
\node[font=\scriptsize, gray, anchor=south] at (axis cs:-1,0.061) {our form};
\addplot[black, line width=0.7pt, mark=*, mark size=1.1pt] table[x=g, y=dgini] {fig13-heavy01.dat};
\addplot[black, line width=0.7pt, mark=o, mark size=1.3pt] table[x=g, y=dgini] {fig13-def1.dat};
\addplot[black, line width=0.7pt, mark=triangle*, mark size=1.4pt] table[x=g, y=dgini] {fig13-even1.dat};
\addplot[black, only marks, mark=*, mark size=1.1pt] coordinates {(-14,0.031)};
\addplot[black, only marks, mark=o, mark size=1.3pt] coordinates {(-14,-0.147)};
\addplot[black, only marks, mark=triangle*, mark size=1.4pt] coordinates {(-14,0.009)};
\node[font=\scriptsize, anchor=north] at (axis cs:-8.5,0.014) {heavy-tailed, tight};
\node[font=\scriptsize, anchor=south] at (axis cs:-10.5,-0.047) {default, ample};
\node[font=\scriptsize, anchor=south] at (axis cs:-7.5,0.042) {evenly spread, ample};
\end{axis}
\end{tikzpicture}
\caption{Complementarity concentrates funding only when the budget is small. The change in the Gini coefficient of the review-informed funder's grants, relative to its Cobb-Douglas level, as the production function moves from Cobb-Douglas ($\gamma = 0$) toward Leontief (detached marks at the left edge), for three field-and-budget combinations. Settings: \ref{app:simulation}.}
\label{fig:concentration}
\end{figure}


% SI Text 6: robustness to the production function (copy of Appendix C; Fig. S5, Table S2).
% Appendix C: robustness to the production technology.
% Referenced by S3, S6, S7, S8. \input inside the appendices environment AFTER
% appendix-lottery. Holds Figure 12 (signal robustness) and Table 2 (timing by
% technology). Grounds: sigma_tierA_gc0 (hash d77c854, 100 seeds), sigma_tierA_gcm3
% (100 seeds), sigma_tierA_leontief (50 seeds), all re-derived in-session
% (verify_s9.R); the tier A sweeps use an earlier version of the allocation routine
% than the body's runs, so magnitudes are compared within this appendix only.
\section{Robustness to the production function}\label{app:technology}

This appendix asks which of our results depend on the harmonic-mean production function of \ref{sec:model}. We embed that function in a family indexed by an exponent $\gamma$: expected output is $A\,M_\gamma(K, R)$, where $M_\gamma(K, R) = \big((K^{\gamma} + R^{\gamma})/2\big)^{1/\gamma}$ is the power mean of capability and resources. At $\gamma = -1$ the power mean is the harmonic mean $2KR/(K + R)$, which is our production function up to the constant $A$. In the limit $\gamma \to 0$ (Cobb-Douglas) it is the geometric mean $\sqrt{KR}$, and a shortfall in either input can be offset by more of the other; as $\gamma$ decreases the inputs grow harder to substitute; in the limit $\gamma \to -\infty$ (Leontief) it is $\min(K, R)$, and output depends only on the scarcer input. We re-run the model's central comparisons across this family.\footnote{Settings: \ref{app:simulation}.}

The value of the review signal survives the whole family (Fig.~\ref{fig:signal-robustness}). At every production function tested, from Cobb-Douglas to Leontief, the signal's value increases as capability grows more heavy-tailed and as review grows more informative. The ordering of review's value by field (main text) is therefore not an artifact of our choice of production function. What changes with the production function is the scale. The more complementary the inputs, the more review is worth: under heavy-tailed capability with reliable review, the signal's value increases from about 16 percent of the research output that funding with review adds at Cobb-Douglas to about 47 percent at $\gamma = -3$ and about half at Leontief.\footnote{15.8, 47.1, and 51.4 percent. Cobb-Douglas lies below the two more complementary production functions at every tested tail and noise level at which the value is distinguishable from zero; $\gamma = -3$ and Leontief are within one standard error of each other at the noisiest review under heavy tails.} Complementarity raises what targeting is worth, and with it the value of the information that targeting requires.

% Figure (Appendix C): the structure of review's value survives the production family.
% y = value of the review signal (S5 - S4) as percent of no-funding output.
% Left panel: vs the capability tail parameter alpha_K at reliable review (tau = 0.3).
% Right panel: vs review noise tau at heavy tails (alpha_K = 1.3).
% Curves: Cobb-Douglas (gamma = 0, filled circles), gamma = -3 (open circles),
% Leontief (triangles). Ground: sigma_tierA_{gc0,gcm3,leontief} signal_value
% summaries (T = 2, greedy allocator n_steps = 400; 100/100/50 seeds), re-derived
% in-session (verify_s9.R). Data: fig12-left.dat, fig12-right.dat.
\begin{figure}[tp]
\centering
\begin{tikzpicture}
\begin{axis}[
  name=left,
  width=0.46\textwidth, height=5.6cm,
  axis lines*=left,
  xmin=1.3, xmax=3.5, ymin=0, ymax=60,
  xtick={1.3, 1.8, 2.5, 3.5},
  ytick={0, 20, 40, 60},
  xlabel={capability tail $\alpha_K$},
  ylabel={value of review (\% of the gain from funding)},
  ylabel style={font=\small}, xlabel style={font=\small},
  tick label style={font=\small},
  axis line style={line width=0.4pt},
  tick style={line width=0.4pt, black},
  title={\small reliable review ($\tau_K = 0.3$)},
  title style={yshift=-2pt},
  clip=false,
]
\addplot[black, line width=0.7pt, mark=*, mark size=1.1pt] table[x=k, y=cd] {fig12-left.dat};
\addplot[black, line width=0.7pt, mark=o, mark size=1.3pt] table[x=k, y=g3] {fig12-left.dat};
\addplot[black, line width=0.7pt, mark=triangle*, mark size=1.4pt] table[x=k, y=leo] {fig12-left.dat};
\node[font=\scriptsize, anchor=south west] at (axis cs:1.33,16.5) {Cobb-Douglas};
\end{axis}
\begin{axis}[
  at={(left.outer east)}, anchor=outer west, xshift=6pt,
  width=0.46\textwidth, height=5.6cm,
  axis lines*=left,
  xmode=log,
  xmin=0.3, xmax=10, ymin=0, ymax=60,
  xtick={0.3, 1, 3, 10},
  xticklabels={0.3, 1, 3, 10},
  ytick={0, 20, 40, 60},
  yticklabels={,,,},
  xlabel={review noise $\tau_K$},
  xlabel style={font=\small},
  tick label style={font=\small},
  axis line style={line width=0.4pt},
  tick style={line width=0.4pt, black},
  title={\small heavy-tailed ($\alpha_K = 1.3$)},
  title style={yshift=-2pt},
  clip=false,
]
\addplot[black, line width=0.7pt, mark=*, mark size=1.1pt] table[x=tau, y=cd] {fig12-right.dat};
\addplot[black, line width=0.7pt, mark=o, mark size=1.3pt] table[x=tau, y=g3] {fig12-right.dat};
\addplot[black, line width=0.7pt, mark=triangle*, mark size=1.4pt] table[x=tau, y=leo] {fig12-right.dat};
\node[font=\scriptsize, anchor=south west] at (axis cs:0.315,16.5) {Cobb-Douglas};
\end{axis}
\end{tikzpicture}
\caption{The value of review across production functions. The value of review as a percent of the research output that review-informed funding adds over no funding, as capability grows more heavy-tailed (left) and as review grows more informative (right), for three production functions: Cobb-Douglas (filled circles), $\gamma = -3$ (open circles), and Leontief (triangles). Settings: \ref{app:simulation}.}
\label{fig:signal-robustness}
\end{figure}


The timing results of the main text do not survive the substitutable end of the family (Table~\ref{tab:timing-technology}). Under Cobb-Douglas production the value of scheduling is indistinguishable from zero at every compounding rate, and the optimal schedule stays within about a hundredth of even. At $\gamma = -3$, a production function more complementary than ours, both the value of scheduling and the shift of spending toward later rounds increase with the compounding rate and exceed their Cobb-Douglas values at every rate. The timing results hold where capability and resources are complements and vanish where they are substitutable.

\begin{table}[t]
\centering
\small
\caption{The value of scheduling and the optimal schedule by production function. Five rounds. The value of scheduling is the forward-looking funder's gain over the myopic funder, both with review, as a percent of the research output that the forward-looking funder's funding adds; the schedule's center of mass is 0.5 when spending is even and larger when spending shifts toward later rounds. One hundred simulated populations per point; standard errors of the gain are at most 0.01 (Cobb-Douglas) and 0.4 ($\gamma = -3$) percentage points.}
\label{tab:timing-technology}
\begin{tabular}{@{}lccccc@{}}
\toprule
compounding rate $\epsilon$ & 0.05 & 0.15 & 0.30 & 0.55 & 0.85 \\
\midrule
\multicolumn{6}{@{}l}{\emph{Value of scheduling (percent of the gain from funding)}} \\
Cobb-Douglas ($\gamma = 0$) & $-0.01$ & $-0.01$ & $0.00$ & $0.01$ & $0.03$ \\
$\gamma = -3$               & $0.05$  & $0.56$  & $1.96$ & $4.56$ & $6.52$ \\
\addlinespace
\multicolumn{6}{@{}l}{\emph{Center of mass of the optimal schedule}} \\
Cobb-Douglas ($\gamma = 0$) & 0.499 & 0.500 & 0.501 & 0.505 & 0.512 \\
$\gamma = -3$               & 0.511 & 0.540 & 0.577 & 0.623 & 0.660 \\
\bottomrule
\end{tabular}
\end{table}

How widely the informed funder spreads its grants across this family is reported in the main text (Fig.~\ref{fig:concentration}).

% SI Text 5: sensitivity of the value of targeting and of review (Tables S4 to S7, Fig. S7).
\section{Sensitivity of the value of targeting and of review}\label{app:sensitivity}\label{si:sensitivity}

The main text varies the tail of the capability distribution, the budget, and review noise, and holds fixed the mean of the capability distribution, the correlation between capability and resources ($\rho = 0$), the tail of the resource distribution ($\alpha_R = 2$), and the size of the applicant pool ($n = 50$). This appendix varies each of these and asks two questions: whether the ordering of fields by the value of targeting and of review survives, and what single quantity, if any, organizes the results. Every run uses two rounds, the default budget $b = 1$, the myopic funders with and without review, and the complete-information benchmark; 50 simulated populations per setting, 100 in the noise sweeps. The value of targeting is the complete-information optimum's gain over uniform funding as a share of uniform funding's gain over no funding; the value of review is the share of the records-only shortfall that the review signal recovers, as in Fig.~2\textit{C}.

\subsection*{Mean capability}
A Pareto distribution with a fixed scale has mean $\alpha/(\alpha - 1)$, so at scale 1 the heavy-tailed field ($\alpha_K = 1.3$) has mean capability 4.3, the default field ($\alpha_K = 2$) 2, and the evenly spread field ($\alpha_K = 3.5$) 1.4: a sweep of $\alpha_K$ at fixed scale varies mean capability along with inequality, and a higher mean makes resources more binding on its own. Table~\ref{tab:sens-mean} compares that sweep with one in which the scale is set to $2(\alpha_K - 1)/\alpha_K$, so that every field has mean capability 2. Holding the mean fixed narrows the heavy-tailed field's advantage but preserves the ordering at every noise level; Fig.~2\textit{C} and every number in the main text use the fixed-mean sweep.

\begin{table}[t]
\centering
\small
\caption{The value of review by field at a fixed Pareto scale and at a fixed mean capability. Share of the records-only shortfall recovered, at review noise $\tau_K = 0.3$, $1$, and $3$; 100 simulated populations per entry.}
\label{tab:sens-mean}
\begin{tabular}{@{}lcccccccc@{}}
\toprule
& \multicolumn{4}{c}{fixed scale 1} & \multicolumn{4}{c}{fixed mean 2} \\
$\alpha_K$ & mean $K$ & $\tau_K = 0.3$ & $1$ & $3$ & mean $K$ & $\tau_K = 0.3$ & $1$ & $3$ \\
\midrule
1.3 & 3.84 & 0.94 & 0.88 & 0.69 & 1.77 & 0.91 & 0.77 & 0.48 \\
2.0 & 1.97 & 0.84 & 0.63 & 0.28 & 1.97 & 0.84 & 0.63 & 0.28 \\
3.5 & 1.39 & 0.56 & 0.21 & 0.03 & 1.99 & 0.62 & 0.31 & 0.05 \\
\bottomrule
\end{tabular}
\end{table}

\subsection*{Review compared at equal informativeness}
The same review noise $\tau_K$ is a more informative signal where capability is more dispersed, because the noise then moves fewer researchers past one another. Table~\ref{tab:sens-matched} therefore compares fields at equal informativeness, measured by the rank correlation between the review score and capability, interpolating within the fixed-mean sweep. At a rank correlation of 0.8, 0.55, and 0.3, the heavy-tailed field recovers 0.91, 0.78, and 0.48 of the shortfall, the default field 0.83, 0.63, and 0.30, and the evenly spread field 0.64, 0.45, and 0.17. The ordering is unchanged; the gaps between fields are smaller than at equal noise.

\begin{table}[t]
\centering
\small
\caption{The value of review by field at equal informativeness. Share of the records-only shortfall recovered when the rank correlation between the review score and capability is held at the stated value; fixed mean capability; interpolated within the noise sweep of Table~\ref{tab:sens-mean}.}
\label{tab:sens-matched}
\begin{tabular}{@{}lccc@{}}
\toprule
$\alpha_K$ & rank correlation 0.8 & 0.55 & 0.3 \\
\midrule
1.3 & 0.91 & 0.78 & 0.48 \\
2.0 & 0.83 & 0.63 & 0.30 \\
3.5 & 0.64 & 0.45 & 0.17 \\
\bottomrule
\end{tabular}
\end{table}

\subsection*{The correlation between capability and resources, the resource tail, and pool size}
Table~\ref{tab:sens-rest} varies the three remaining parameters one at a time, at the default review noise and at Pareto scale 1. The correlation between capability and resources is the one that matters. As $\rho$ rises from $-0.5$ to $0.8$, the value of targeting falls in every field (from 0.96 to 0.69 in the heavy-tailed field, from 0.25 to 0.06 in the evenly spread one), and so does the value of review; the ordering of fields is unchanged at every $\rho$. The mechanism is the gap rule itself: when resources already sit with the capable, the gaps $cK_i - R_{i0}$ are small and even, and there is little for targeting, or for the information that targeting requires, to add. The tail of the resource distribution changes the value of targeting by at most a tenth and the value of review by less. Pool size raises the value of targeting in the heavy-tailed field, because a larger pool contains more of the rare researchers of very high capability, and raises the value of review in the evenly spread field, where a larger pool gives review more to separate; neither effect changes the ordering.

\begin{table}[t]
\centering
\small
\caption{The value of targeting (upper block) and of review (lower block) by field as the correlation between capability and resources, the tail of the resource distribution, and the size of the applicant pool vary. Review noise $\tau_K = 1$; Pareto scale 1; 50 simulated populations per entry. Default values: $\rho = 0$, $\alpha_R = 2$, $n = 50$.}
\label{tab:sens-rest}
\begin{tabular}{@{}lcccccccccccc@{}}
\toprule
& \multicolumn{4}{c}{correlation $\rho$} & \multicolumn{3}{c}{resource tail $\alpha_R$} & \multicolumn{4}{c}{pool size $n$} \\
$\alpha_K$ & $-0.5$ & $0$ & $0.5$ & $0.8$ & $1.3$ & $2$ & $3.5$ & $25$ & $50$ & $100$ & $200$ \\
\midrule
\multicolumn{12}{@{}l}{\emph{Value of targeting (share of uniform funding's gain)}} \\
1.3 & 0.96 & 0.91 & 0.82 & 0.69 & 0.92 & 0.91 & 0.82 & 0.81 & 0.91 & 1.06 & 1.15 \\
2.0 & 0.56 & 0.49 & 0.38 & 0.25 & 0.43 & 0.49 & 0.46 & 0.44 & 0.49 & 0.53 & 0.57 \\
3.5 & 0.25 & 0.20 & 0.13 & 0.06 & 0.17 & 0.20 & 0.18 & 0.20 & 0.20 & 0.21 & 0.22 \\
\addlinespace
\multicolumn{12}{@{}l}{\emph{Value of review (share of the records-only shortfall recovered)}} \\
1.3 & 0.89 & 0.89 & 0.86 & 0.84 & 0.86 & 0.89 & 0.91 & 0.86 & 0.89 & 0.90 & 0.89 \\
2.0 & 0.64 & 0.63 & 0.59 & 0.54 & 0.64 & 0.63 & 0.64 & 0.58 & 0.63 & 0.68 & 0.69 \\
3.5 & 0.22 & 0.20 & 0.14 & 0.00 & 0.22 & 0.20 & 0.19 & 0.14 & 0.20 & 0.25 & 0.30 \\
\bottomrule
\end{tabular}
\end{table}

\subsection*{Review that partly reflects resources}
The main text's review signal observes capability with noise. Real review may also reflect the resources a researcher visibly has. Table~\ref{tab:sens-contam} replaces the signal by $K_i + \beta R_{i0} + \xi_{it}$, with the funder knowing $\beta$ and updating accordingly. Contamination costs little where capability is heavy-tailed: at $\beta = 1$, where the score weights resources as heavily as capability, the heavy-tailed field still recovers 0.60 to 0.65 of the shortfall, against 0.77 to 0.91 for a clean signal. It costs most where capability is evenly spread, where at $\beta = 1$ review is nearly worthless. The reason is the one behind Fig.~2\textit{C}: the researchers whom review must find in a heavy-tailed field stand so far apart in capability that a resource term does not hide them.

\begin{table}[t]
\centering
\small
\caption{The value of review when the review score partly reflects resources. Share of the records-only shortfall recovered when the score is $K_i + \beta R_{i0}$ plus noise, for $\beta = 0$ (the main text's signal), $0.25$, $0.5$, and $1$; fixed mean capability; 50 simulated populations per entry (100 at $\beta = 0$).}
\label{tab:sens-contam}
\begin{tabular}{@{}lcccccccc@{}}
\toprule
& \multicolumn{4}{c}{$\tau_K = 0.3$} & \multicolumn{4}{c}{$\tau_K = 1$} \\
$\alpha_K$ & $\beta = 0$ & $0.25$ & $0.5$ & $1$ & $\beta = 0$ & $0.25$ & $0.5$ & $1$ \\
\midrule
1.3 & 0.91 & 0.89 & 0.81 & 0.60 & 0.77 & 0.76 & 0.73 & 0.65 \\
2.0 & 0.84 & 0.79 & 0.66 & 0.38 & 0.63 & 0.61 & 0.57 & 0.45 \\
3.5 & 0.62 & 0.52 & 0.33 & 0.04 & 0.31 & 0.26 & 0.22 & 0.11 \\
\bottomrule
\end{tabular}
\end{table}

\subsection*{The inequality of the gaps organizes the results}
The runs above vary five parameters. One quantity summarizes what they do to the value of targeting: how unequally the optimal grants, the gaps $\max(cK_i - R_{i0}, 0)$ at the population's budget, are distributed across the field. Across all 159 settings in this appendix, the logarithm of the value of targeting is a nearly linear function of the Gini coefficient of the optimal grants, with $R^2 = 0.96$ (Fig.~\ref{fig:gap-gini}); adding the ratio of mean capability to mean resources raises $R^2$ to $0.99$. Settings that differ in the capability tail, the correlation between capability and resources, the resource tail, and pool size, but agree on the inequality of the gaps, have the same value of targeting. The value of review depends on the same quantity together with the signal's informativeness: the log-odds of the share recovered, regressed on the Gini coefficient of the optimal grants and the log-odds of the rank correlation between score and capability, has $R^2 = 0.86$, and adding the capability tail as a separate variable raises it by less than $0.01$. Capability inequality is the main driver of gap inequality in every sweep we ran, which is why the main text can speak of it directly; the correlation between capability and resources is the second driver, and it works through the same channel.

% Figure (Appendix E): the value of targeting against the Gini coefficient of the optimal
% grants, across the 159 settings of Appendix E (capability tail, capability-resource
% correlation, resource tail, pool size, mean normalization). Marker by capability tail.
% Data: fig15-gap-gini.dat (gini, targ, share, skcor, a) from sens.R.
\begin{figure}[tp]
\centering
\pgfplotsset{discard if not/.style 2 args={x filter/.append code={\edef\tempa{\thisrow{#1}}\edef\tempb{#2}\ifx\tempa\tempb\else\def\pgfmathresult{inf}\fi}}}
\begin{tikzpicture}
\begin{axis}[
  width=0.78\textwidth, height=6.0cm,
  axis lines*=left,
  ymode=log,
  xmin=0.25, xmax=0.95, ymin=0.04, ymax=2,
  xtick={0.3, 0.5, 0.7, 0.9},
  ytick={0.05, 0.1, 0.2, 0.5, 1, 2},
  yticklabels={0.05, 0.1, 0.2, 0.5, 1, 2},
  xlabel={Gini coefficient of the optimal grants},
  ylabel={value of targeting (share of uniform funding's gain)},
  ylabel style={font=\small}, xlabel style={font=\small},
  tick label style={font=\small},
  axis line style={line width=0.4pt},
  tick style={line width=0.4pt, black},
  clip=false,
]
\addplot[only marks, mark=*, mark size=1.4pt, black, discard if not={a}{1.3}] table[x=gini, y=targ] {fig15-gap-gini.dat};
\addplot[only marks, mark=o, mark size=1.5pt, black, line width=0.5pt, discard if not={a}{2}] table[x=gini, y=targ] {fig15-gap-gini.dat};
\addplot[only marks, mark=triangle*, mark size=1.6pt, gray, discard if not={a}{3.5}] table[x=gini, y=targ] {fig15-gap-gini.dat};
\node[font=\small, anchor=west] at (axis cs:0.27,1.55) {$\alpha_K = 1.3$ (filled)};
\node[font=\small, anchor=west] at (axis cs:0.27,1.05) {$\alpha_K = 2$ (open)};
\node[font=\small, anchor=west] at (axis cs:0.27,0.72) {$\alpha_K = 3.5$ (triangles)};
\end{axis}
\end{tikzpicture}
\caption{The value of targeting against the inequality of the optimal grants. Each point is one of the 159 settings of \ref{app:sensitivity}: the value of targeting (logarithmic scale) against the Gini coefficient of the optimal grants at that setting's budget. The settings vary the capability tail (marker), the association between capability and resources, the resource tail, the size of the applicant pool, and the scale of capability. The logarithm of the value of targeting is linear in the Gini coefficient with $R^2 = 0.96$.}
\label{fig:gap-gini}
\end{figure}


% SI Text 6: strategic timing of funding. Moved from the main text's Results on 2026-09-29
% (revision plan 2, Package B). The text is the former Results subsection with three stated
% limits added; the figure is the former main-text Fig. 4. Terms: "the forward-looking
% funder's timing" (not "optimal timing"); ``gain from planning ahead'' (not "timing gain").
\section{Strategic timing of funding}\label{si:timing}
The funder holds a fixed budget for the whole horizon and may time its spending freely. The strongest argument for spending early concerns fields whose researchers have few resources: a researcher without resources produces nothing, so the funder observes nothing and no capability grows, and money held for later rounds produces neither evidence nor growth in the meantime. The argument fails. A funder that spends everything in the first round allocates before observing any research output, and knows exactly as little as a funder that spends everything in the last round; in a field with no resources at all, the two produce identical expected research output. What the argument does establish is a complementarity: early spending enables observation and growth, and late spending exploits what early spending revealed. That favors putting some money early. It does not favor putting most of it early.

In our simulations, under complementarity, observation comes first and money second. Wherever capabilities compound at more than a negligible rate ($\epsilon$ above about 0.03), the forward-looking funder spends less than an even share in the first round and spends the bulk once research output has revealed who is capable; scarce baseline resources raise the early share and weaken the effect, but do not reverse it (Fig.~\ref{fig:p3}\textit{A}). The shift of spending toward later rounds comes from the research output researchers produce whether or not they are funded, which compounds and reveals capability while the funder waits (\ref{si:extended}, Fig.~\ref{fig:depth-schedule}).

Planning ahead is worth less than the signal provided by grant peer review, in every setting we test. Across all horizons and compounding rates in our sweeps, the forward-looking funder's gain over the myopic funder, which spends equal installments each round, stays below the review signal's value: by an order of magnitude or more at the default compounding rate and below, and approaching the signal's value only at the highest compounding rate and the longest horizon we consider (Fig.~\ref{fig:p3}\textit{B}). As a share of the research output that the forward-looking funder's funding adds, the gain from planning ahead ranges from indistinguishable from zero to 4.5\% across the 30 settings tested; the review signal's value in the same settings ranges from 7\% to 27\%. The gain exists only where capability and resources are complements: under Cobb-Douglas production it is indistinguishable from zero (\ref{app:technology}, Table~\ref{tab:timing-technology}).

Three limits bound these results, and they are the reason the timing analysis sits in the SI rather than the main text. First, the forward-looking funder is a certainty-equivalent planner (SI Materials and Methods): it treats future research output as equal to its expectation and does not value the information its own grants will generate, so it is not the Bayes-optimal sequential policy, and its gain is a lower bound on what planning ahead could be worth. Second, the gain combines two changes at once, when the money is spent and whom it reaches, because a funder that anticipates capability growth also values recipients differently; we do not separate the two, and ``gain from planning ahead'' names their sum. Third, the objective counts research output within the horizon only and assigns no value to the capability that remains at its end. A terminal value for capability would credit the capability that late grants build, which the horizon objective ignores, and also the compounded capability that early grants leave behind; we have not measured which effect is larger, so the size of the shift toward later rounds should be read as specific to the horizon objective.

\begin{figure}[tp]
\centering
\begin{minipage}{\textwidth}\noindent\textbf{A}\par\vspace{-0.3em}
\begin{tikzpicture}
\begin{axis}[
  width=0.79\textwidth, height=5.4cm,
  axis lines*=left,
  xmode=log,
  xmin=0.0001, xmax=0.85, ymin=0.48, ymax=0.68,
  xtick={0.0001, 0.01, 0.1, 0.85},
  xticklabels={0.0001, 0.01, 0.1, 0.85},
  ytick={0.5, 0.55, 0.6, 0.65},
  yticklabels={even, 0.55, 0.60, 0.65},
  xlabel={compounding rate $\epsilon$},
  ylabel={center of mass of spending},
  ylabel style={font=\small}, xlabel style={font=\small},
  tick label style={font=\small},
  axis line style={line width=0.4pt},
  tick style={line width=0.4pt, black},
  clip=false,
]
\addplot[gray, dashed, line width=0.4pt, domain=0.0001:0.85] {0.5};
\addplot[black, line width=0.7pt] table[x=eps, y=rich] {fig8-boundary.dat};
\addplot[black, line width=0.7pt] table[x=eps, y=mid]  {fig8-boundary.dat};
\addplot[black, line width=0.7pt] table[x=eps, y=poor] {fig8-boundary.dat};
\node[font=\small, anchor=west] at (axis cs:0.92,0.652) {many resources};
\node[font=\small, anchor=west] at (axis cs:0.92,0.566) {intermediate};
\node[font=\small, anchor=west] at (axis cs:0.92,0.521) {few resources};
\end{axis}
\end{tikzpicture}
\end{minipage}\par\vspace{0.9em}
\begin{minipage}{\textwidth}\noindent\textbf{B}\par\vspace{-0.3em}
\begin{tikzpicture}
\begin{axis}[
  width=0.82\textwidth, height=5.7cm,
  axis lines*=left,
  xmin=2, xmax=10, ymin=0, ymax=1.05,
  xtick={2, 4, 6, 8, 10},
  ytick={0, 0.25, 0.5, 0.75, 1},
  yticklabels={0, 0.25, 0.5, 0.75, 1},
  xlabel={horizon $T$ (rounds)},
  ylabel={gain from planning ahead / signal value},
  ylabel style={font=\small}, xlabel style={font=\small},
  tick label style={font=\small},
  axis line style={line width=0.4pt},
  tick style={line width=0.4pt, black},
  clip=false,
]
\addplot[gray, dashed, line width=0.4pt, domain=2:10] {1};
\node[font=\small, gray, anchor=south west] at (axis cs:2.1,1.005) {equal to the signal's value};
\addplot[black, line width=0.7pt] coordinates {(2,0.0468) (3,0.1136) (4,0.1855) (5,0.2635) (6,0.3443) (7,0.4213) (8,0.4982) (9,0.5811) (10,0.6556)};
\addplot[black, line width=0.7pt] coordinates {(2,0.0198) (3,0.0482) (4,0.0769) (5,0.1090)};
\addplot[black, line width=0.7pt] coordinates {(2,0.0058) (3,0.0144) (4,0.0224) (5,0.0317) (6,0.0398) (7,0.0487) (8,0.0561) (9,0.0651) (10,0.0749)};
\addplot[black, line width=0.7pt] coordinates {(2,0.0012) (3,0.0034) (4,0.0051) (5,0.0069)};
\addplot[black, line width=0.7pt] coordinates {(2,0.0000) (3,0.0002) (4,0.0001) (5,0.0000)};
\node[font=\small, anchor=west] at (axis cs:10.2,0.656) {$\epsilon = 0.85$};
\node[font=\small, anchor=west] at (axis cs:5.15,0.126) {$\epsilon = 0.55$};
\node[font=\small, anchor=west] at (axis cs:10.2,0.075) {$\epsilon = 0.3$};
\end{axis}
\end{tikzpicture}
\end{minipage}\par\vspace{0.9em}
\caption{Strategic timing of funding. (\textit{A}) The compounding rate, not the scarcity of baseline resources, sets how far spending shifts toward later rounds: the center of mass of spending (0.5 is even; higher is later) as the compounding rate $\epsilon$ increases, for three baseline resource scales. (\textit{B}) Choosing whom outweighs choosing when: the gain from planning ahead, as a share of the review signal's value, by horizon and compounding rate; in every setting the gain is below the signal's value. Settings: \ref{app:simulation}.}
\label{fig:p3}
\end{figure}



\bibliographystyle{plainnat}
\bibliography{references}
\end{document}
